\section{Class PSPACE --- TQBF}
%%% generated by the blueprint pipeline from TCSlib.Complexity.ClassPSPACE.TQBF
\subsection*{Class PSPACE --- TQBF: Introduction}

Declarations of the Lean module \texttt{TCSlib.Complexity.ClassPSPACE.TQBF}.

\subsection*{Class PSPACE --- TQBF: Results}

\begin{definition}[PSPACE hardness]
\lean{Complexity.PSPACEHard}
\label{Complexity.PSPACEHard}
\leanok
\uses{Complexity.PSPACE}
A language $L'$ over the Boolean alphabet is PSPACE-hard if every language in PSPACE has
a polynomial-time many-one reduction to $L'$.
\end{definition}

\begin{definition}[PSPACE completeness]
\lean{Complexity.PSPACEComplete}
\label{Complexity.PSPACEComplete}
\leanok
\uses{Complexity.PSPACE, Complexity.PSPACEHard}
A language $L'$ over the Boolean alphabet is PSPACE-complete if it belongs to PSPACE and
every language in PSPACE has a polynomial-time many-one reduction to $L'$.
\end{definition}

\begin{theorem}[A polynomial-space complete language in P]
\lean{Complexity.PSPACE_eq_P_of_pspaceComplete_mem_P}
\label{Complexity.PSPACE_eq_P_of_pspaceComplete_mem_P}
\leanok
\uses{Complexity.P, Complexity.PSPACE, Complexity.PSPACEComplete, Complexity.P_subset_PSPACE, Complexity.mem_P_of_polyTimeReducible}
\difficulty{3}
Let $L'$ be a language over the binary alphabet that is polynomial-space complete: $L'$
is decidable in polynomial space, and every language decidable in polynomial space has a
polynomial-time many-one reduction to $L'$. If $L'$ is also decidable in polynomial
time, then the class of polynomial-space decidable languages equals the class of
polynomial-time decidable languages.
\end{theorem}

\begin{definition}[Machine state count]
\lean{Complexity.ConfigCodec.stateCount}
\label{Complexity.ConfigCodec.stateCount}
\leanok
For a finite-state Turing machine $M$, define its state count to be the number of states
of $M$ plus one.
\end{definition}

\begin{definition}[Machine state index]
\lean{Complexity.ConfigCodec.stateIndex}
\label{Complexity.ConfigCodec.stateIndex}
\leanok
For a finite-state Turing machine $M$, assign index $0$ to the absent state. Assign each
actual state the index of its position in a finite enumeration of the states, plus one.
\end{definition}

\begin{definition}[State offset]
\lean{Complexity.ConfigCodec.stateOffset}
\label{Complexity.ConfigCodec.stateOffset}
\leanok
For a finite-state Turing machine $M$ and a natural number $n$, define the state offset
to be $2n+2$.
\end{definition}

\begin{definition}[Work offset]
\lean{Complexity.ConfigCodec.workOffset}
\label{Complexity.ConfigCodec.workOffset}
\leanok
\uses{Complexity.ConfigCodec.stateCount, Complexity.ConfigCodec.stateOffset}
For a finite-state Turing machine $M$ and a natural number $n$, define the work offset
to be $2n+2$ plus the number of states of $M$ plus one.
\end{definition}

\begin{definition}[Summary offset]
\lean{Complexity.ConfigCodec.summaryOffset}
\label{Complexity.ConfigCodec.summaryOffset}
\leanok
\uses{Complexity.ConfigCodec.workOffset}
Let $M$ be a Turing machine with Boolean alphabet, and let $n,s\in\mathbb{N}$. Define
the summary offset to be the work offset of $M$ at $n$, plus three times the product of
the machine parameter $k$ and $2s+1$.
\end{definition}

\begin{definition}[Configuration encoding width]
\lean{Complexity.ConfigCodec.width}
\label{Complexity.ConfigCodec.width}
\leanok
\uses{Complexity.ConfigCodec.summaryOffset}
Let $M$ be a Turing machine with Boolean symbols, and let $n,s\in\mathbb{N}$. Define the
width associated with $M,n,s$ to be the summary offset for $M,n,s$, increased by $2$.
\end{definition}

\begin{definition}[Input variable numbering]
\lean{Complexity.ConfigCodec.inputVar}
\label{Complexity.ConfigCodec.inputVar}
\leanok
For each natural number $i$, define the input-variable number to be $i$.
\end{definition}

\begin{definition}[Head variable number]
\lean{Complexity.ConfigCodec.headVar}
\label{Complexity.ConfigCodec.headVar}
\leanok
For natural numbers $n$ and $j$, define the head variable number to be $n+j$.
\end{definition}

\begin{definition}[State variable numbering]
\lean{Complexity.ConfigCodec.stateVar}
\label{Complexity.ConfigCodec.stateVar}
\leanok
\uses{Complexity.ConfigCodec.stateIndex, Complexity.ConfigCodec.stateOffset}
For a natural number $n$ and an optional machine state $q$, define its state-variable
number to be the state-block offset for input length $n$ plus the index assigned to $q$.
\end{definition}

\begin{definition}[Work-tape variable numbering]
\lean{Complexity.ConfigCodec.workVar}
\label{Complexity.ConfigCodec.workVar}
\leanok
\uses{Complexity.ConfigCodec.workOffset}
Let $M$ be a deterministic Turing machine with Boolean-valued tape symbols, and let
$n,s,i,j,b$ be natural numbers. The work-variable number is
$\mathrm{workOffset}(M,n)+3((2s+1)i+j)+b$, where $\mathrm{workOffset}(M,n)$ is the
work-block offset for $M$ and input length $n$.
\end{definition}

\begin{definition}[Summary variable numbering]
\lean{Complexity.ConfigCodec.summaryVar}
\label{Complexity.ConfigCodec.summaryVar}
\leanok
\uses{Complexity.ConfigCodec.summaryOffset}
For natural numbers $n,s$, define the summary-variable number for a Boolean flag $d$ to
be the summary-block offset for the machine, input length $n$, and parameter $s$, plus
$1$ when $d$ is true and $0$ when it is false.
\end{definition}

\begin{definition}[Work-tape cell flags]
\lean{Complexity.ConfigCodec.workFlags}
\label{Complexity.ConfigCodec.workFlags}
\leanok
\uses{Turing.Cfg}
Given natural numbers $s,j$, a configuration of a $k$-tape machine on a Boolean input,
and a work-tape index $i$, associate to the cell at position $j-s$ the three Boolean
flags indicating, respectively, that the cell contains false, that it contains true, and
that the work-tape head is at that position.
\end{definition}

\begin{definition}[Configuration encoding]
\lean{Complexity.ConfigCodec.code}
\label{Complexity.ConfigCodec.code}
\leanok
\uses{Complexity.ConfigCodec.stateCount, Complexity.ConfigCodec.stateIndex, Complexity.ConfigCodec.workFlags, Turing.Cfg, Turing.outSummary}
For a machine, natural numbers $n,s$, a Boolean list $x$, and a configuration $c$ on
input $x$, define its Boolean code by concatenating: the first $n$ input bits, using
false for positions beyond the end of $x$; the one-hot encoding of the input-head
position among $n+2$ positions; the one-hot encoding of the state, including the
optional halting state; the three work-tape flags for each of the machine's tapes and
each position indexed from $0$ through $2s$; and the two flags indicating whether the
output summary is acceptance or deadness.
\end{definition}

\begin{definition}[Configuration validity clauses]
\lean{Complexity.ConfigCodec.validity}
\label{Complexity.ConfigCodec.validity}
\leanok
\uses{Complexity.ConfigCodec.headVar, Complexity.ConfigCodec.stateCount, Complexity.ConfigCodec.stateOffset, Complexity.ConfigCodec.summaryVar, Complexity.ConfigCodec.workVar, Complexity.FormulaKit.exactlyOne}
For a machine and natural numbers $n,s$, define a CNF formula by concatenating clauses
that require exactly one of the $n+2$ input-head positions and exactly one of the
machine's state indices to be selected; for each work tape, exactly one of the $2s+1$
head-position flags to be selected; for every work-tape cell in that range, prohibit
either of its two symbol flags from being true; and require at least one of the two
summary flags to be true.
\end{definition}

\begin{definition}[Acceptance clauses]
\lean{Complexity.ConfigCodec.acceptance}
\label{Complexity.ConfigCodec.acceptance}
\leanok
\uses{Complexity.ConfigCodec.stateOffset, Complexity.ConfigCodec.summaryVar}
For a Boolean-symbol Turing machine $M$ and natural numbers $n,s$, let $a(n)$ denote the
state variable offset and let $b_M(n,s,0)$ denote the summary variable associated with
the false value. The acceptance formula is the conjunction of the two unit clauses
requiring both variables to be true.
\end{definition}

\begin{definition}[Input equality clauses]
\lean{Complexity.ConfigCodec.inputEquality}
\label{Complexity.ConfigCodec.inputEquality}
\leanok
\uses{Complexity.ConfigCodec.width}
For a Boolean-symbol Turing machine $M$ and natural numbers $n,s$, let $w_M(n,s)$ denote
the width used to encode a configuration. For each $i$ with $0\le i<n$, include the two
clauses $(\neg x_i\lor y_i)$ and $(x_i\lor\neg y_i)$, where $y_i$ is the variable with
index $w_M(n,s)+i$; their conjunction enforces equality of the two bits at each such
position.
\end{definition}

\begin{definition}[Machine configuration cases]
\lean{Complexity.ConfigCodec.Case}
\label{Complexity.ConfigCodec.Case}
\leanok
For a machine $M$, a case consists of an optional state of $M$, an optional input symbol
in the Boolean alphabet, and, for each of its $k$ work tapes, an optional Boolean
symbol.
\end{definition}

\begin{definition}[Action associated to a case]
\lean{Complexity.ConfigCodec.caseAction}
\label{Complexity.ConfigCodec.caseAction}
\leanok
\uses{Complexity.ConfigCodec.Case, Turing.Action}
For a machine $M$ and a case $q$, associate the machine action for the state and scanned
symbols recorded by $q$ when its state is present. If its state is absent, the
associated action has zero input-head movement, zero movement on every work tape, no
writes, no output, and no successor state.
\end{definition}

\begin{definition}[List of all machine cases]
\lean{Complexity.ConfigCodec.caseList}
\label{Complexity.ConfigCodec.caseList}
\leanok
\uses{Complexity.ConfigCodec.Case}
For a machine $M$, form a list containing every possible case exactly once, in the order
induced by listing the elements of the finite case space.
\end{definition}

\begin{definition}[Configuration encoding ports]
\lean{Complexity.ConfigCodec.Port}
\label{Complexity.ConfigCodec.Port}
\leanok
For a natural number $k$, define the ports used to encode a machine configuration: a
state, the input-head position and scanned input symbol, a head position for each work
tape, a port for each work-tape index and Boolean symbol, a summary port for each
Boolean value, and old-bit and new-bit ports.
\end{definition}

\begin{definition}[Number of encoding ports]
\lean{Complexity.ConfigCodec.portCount}
\label{Complexity.ConfigCodec.portCount}
\leanok
For a natural number $k$, the port count is $3k+7$.
\end{definition}

\begin{definition}[Indices of encoding ports]
\lean{Complexity.ConfigCodec.portIndex}
\label{Complexity.ConfigCodec.portIndex}
\leanok
\uses{Complexity.ConfigCodec.Port}
For a natural number $k$ and a port, assign indices as follows: state, input head, and
input symbol receive $0,1,2$; work-head port $i$ receives $3+i$; work-symbol port
$(i,t)$ receives $3+k+2i+t$; summary port $t$ receives $3+3k+t$; and the old-bit and
new-bit ports receive $5+3k$ and $6+3k$, respectively, where each Boolean value is
identified with $0$ or $1$.
\end{definition}

\begin{definition}[Configuration port encoding]
\lean{Complexity.ConfigCodec.portOf}
\label{Complexity.ConfigCodec.portOf}
\leanok
\uses{Complexity.ConfigCodec.Port}
For a machine with $k$ work tapes, indices $0,1,2$ denote the state, input-head, and input-symbol ports; indices $3$ through $2+k$ denote work-head ports. Indices $j$ with $3+k\le j<3+3k$ denote work-symbol ports: the tape number is $\left\lfloor(j-3-k)/2\right\rfloor$, and the Boolean label is true exactly when $j-3-k$ is odd. The next two indices denote the false and true summary ports, followed by the old-bit and new-bit ports; all larger indices also denote the new-bit port.
\end{definition}

\begin{definition}[Expected output categories]
\lean{Complexity.ConfigCodec.Expect}
\label{Complexity.ConfigCodec.Expect}
\leanok
An expectation is one of three kinds: copying a bit, requiring a specified constant
Boolean value, or requiring a summary category indexed by a Boolean flag.
\end{definition}

\begin{definition}[Gate kind data]
\lean{Complexity.ConfigCodec.GateKind}
\label{Complexity.ConfigCodec.GateKind}
\leanok
\uses{Complexity.ConfigCodec.Case, Complexity.ConfigCodec.Expect}
For a machine $M$, a gate kind consists of a pair of optional state, input-symbol, and
work-tape-symbol data describing a partial configuration, a Boolean value, and an
expected output category.
\end{definition}

\begin{definition}[Reading a Boolean symbol]
\lean{Complexity.ConfigCodec.readSymbol}
\label{Complexity.ConfigCodec.readSymbol}
\leanok
Given two Boolean values $f$ and $t$, return true if $t$ is true; otherwise return false
if $f$ is true; if both are false, return no value.
\end{definition}

\begin{definition}[Reading an output summary]
\lean{Complexity.ConfigCodec.readSummary}
\label{Complexity.ConfigCodec.readSummary}
\leanok
\uses{Turing.OutSummary}
Given Boolean flags for acceptance and deadness, return the dead category if the
deadness flag is true; otherwise return the acceptance category if the acceptance flag
is true; otherwise return the empty category.
\end{definition}

\begin{definition}[Summary word encoding]
\lean{Complexity.ConfigCodec.summaryWord}
\label{Complexity.ConfigCodec.summaryWord}
\leanok
\uses{Turing.OutSummary}
Encode the empty output summary by the empty Boolean word, the acceptance summary by the
one-letter word consisting of true, and the dead summary by the one-letter word
consisting of false.
\end{definition}

\begin{definition}[Configuration guard]
\lean{Complexity.ConfigCodec.guard}
\label{Complexity.ConfigCodec.guard}
\leanok
\uses{Complexity.ConfigCodec.Case, Complexity.ConfigCodec.Port, Complexity.ConfigCodec.readSymbol}
Let $q$ be a partial configuration of a machine with $k$ work tapes, let $interior$ be a
Boolean flag, and assign a Boolean value to each configuration port. The guard is true
exactly when the state, input-head, and every work-head port are true; the input-symbol
port, when $interior$ is true, agrees with the optional input symbol specified by $q$
(and when $interior$ is false, that optional symbol is absent); and, for each work tape,
reading its two Boolean work-symbol ports yields the optional symbol specified by $q$.
\end{definition}

\begin{definition}[Expected bit predicate]
\lean{Complexity.ConfigCodec.expectBit}
\label{Complexity.ConfigCodec.expectBit}
\leanok
\uses{Complexity.ConfigCodec.Expect, Complexity.ConfigCodec.Port, Complexity.ConfigCodec.readSummary, Complexity.ConfigCodec.summaryWord, Turing.Action, Turing.outSummary}
Given a machine action and a Boolean assignment to the ports, an expectation holds as
follows: a copy expectation reads the old-bit port; a constant expectation returns its
specified Boolean value; and a summary expectation reads the current summary flags,
encodes the resulting summary as a word, appends the action's optional output symbol if
present, and returns whether the summary of that word is dead when the flag is true and
accepting otherwise.
\end{definition}

\begin{definition}[Transition-kind predicate]
\lean{Complexity.ConfigCodec.kindPred}
\label{Complexity.ConfigCodec.kindPred}
\leanok
\uses{Complexity.ConfigCodec.GateKind, Complexity.ConfigCodec.Port, Complexity.ConfigCodec.caseAction, Complexity.ConfigCodec.expectBit, Complexity.ConfigCodec.guard}
For a finite Turing machine $M$, a gate kind consists of a machine case, a Boolean flag,
and an expectation. On an assignment of Boolean values to the ports, the predicate is
true exactly when the guard for that case and flag is false, or the value at the new-bit
port equals the bit prescribed by the case’s machine action and expectation.
\end{definition}

\begin{definition}[Base-(2s+1) quotient remainder]
\lean{Complexity.ConfigCodec.headOf}
\label{Complexity.ConfigCodec.headOf}
\leanok
For natural numbers $s$, $h$, and $i$, define $\mathrm{headOf}(s,h,i)$ to be the
remainder of $\left\lfloor h/(2s+1)^i\right\rfloor$ upon division by $2s+1$.
\end{definition}

\begin{definition}[Clamped input-head move]
\lean{Complexity.ConfigCodec.moveInput}
\label{Complexity.ConfigCodec.moveInput}
\leanok
For natural numbers $n,p$ and a sign $m$, first add the integer represented by $m$ to
$p$ and truncate a negative result to zero. The resulting position is returned if it is
less than $n+2$; otherwise the returned position is $n+1$.
\end{definition}

\begin{definition}[Expected bit at a wire]
\lean{Complexity.ConfigCodec.expectOf}
\label{Complexity.ConfigCodec.expectOf}
\leanok
\uses{Complexity.ConfigCodec.Case, Complexity.ConfigCodec.Expect, Complexity.ConfigCodec.caseAction, Complexity.ConfigCodec.headOf, Complexity.ConfigCodec.moveInput, Complexity.ConfigCodec.stateIndex, Complexity.ConfigCodec.stateOffset, Complexity.ConfigCodec.summaryOffset, Complexity.ConfigCodec.workOffset}
For a finite Turing machine $M$, natural numbers $n,s,p,h,v$, and a machine case $q$,
let the action be the one prescribed by $q$. The expectation at wire $v$ is a copy when
$v<n$. For $n\le v<2n+2$, it is the constant asserting that the moved input-head
position equals $v-n$; for $2n+2\le v<2n+2+\lvert M.\mathrm{State}\rvert+1$, it is the
constant asserting that the action’s state index equals the offset of $v$ from $2n+2$.
In the work-tape region, decode the tape number, cell, and one of three wire types from
the offset of $v$: out-of-range tape numbers are copied; for the first two wire types,
cells other than the encoded head position are copied, and at that position the result
is copied for an unchanged symbol or is the constant specifying the action’s written
symbol; for the third wire type, the result is the constant asserting that the head
position after the action’s signed move equals the decoded cell. At every remaining
wire, the expectation is the summary expectation whose flag is true exactly when $v$ is
one more than the summary-region offset.
\end{definition}

\begin{definition}[Wire assigned to each port]
\lean{Complexity.ConfigCodec.portWire}
\label{Complexity.ConfigCodec.portWire}
\leanok
\uses{Complexity.ConfigCodec.Case, Complexity.ConfigCodec.Port, Complexity.ConfigCodec.headOf, Complexity.ConfigCodec.headVar, Complexity.ConfigCodec.inputVar, Complexity.ConfigCodec.stateVar, Complexity.ConfigCodec.summaryVar, Complexity.ConfigCodec.width, Complexity.ConfigCodec.workVar}
For a finite Turing machine $M$, natural numbers $n,s,p,h,v$, and a machine case $q$,
assign wires to ports as follows. The state port uses the wire for the state in $q$; the
input-head port uses the wire for position $p$; and the input-symbol port uses the wire
for input cell $p-1$ when $0<p\le n$, otherwise the non-dead summary wire. A work-head
port for tape $i$ uses the wire for the encoded head cell of tape $i$ and its head
marker; a work-symbol port for tape $i$ and symbol $t$ uses the corresponding symbol
wire at that cell. A summary port uses the wire for its Boolean flag, the old-bit port
uses wire $v$, and the new-bit port uses wire $v$ shifted by the total width.
\end{definition}

\begin{definition}[CNF encoding of a transition gate]
\lean{Complexity.ConfigCodec.transitionGate}
\label{Complexity.ConfigCodec.transitionGate}
\leanok
\uses{Complexity.ConfigCodec.Case, Complexity.ConfigCodec.expectOf, Complexity.ConfigCodec.kindPred, Complexity.ConfigCodec.portCount, Complexity.ConfigCodec.portIndex, Complexity.ConfigCodec.portOf, Complexity.ConfigCodec.portWire, Complexity.FormulaKit.gate, Complexity.FormulaKit.readBit}
For a finite Turing machine $M$, natural numbers $n,s,p,h,v$, and a machine case $q$,
form the CNF gate encoding of the Boolean predicate on port assignments that accepts
when the transition-kind predicate is true. Read each gate-input bit from the assignment
at the wire corresponding to that port, and connect each port to its assigned wire.
\end{definition}

\begin{definition}[CNF encoding of all transitions]
\lean{Complexity.ConfigCodec.transition}
\label{Complexity.ConfigCodec.transition}
\leanok
\uses{Complexity.ConfigCodec.caseList, Complexity.ConfigCodec.transitionGate, Complexity.ConfigCodec.width}
For a finite Turing machine $M$ and natural numbers $n,s$, concatenate the
transition-gate CNF encodings over every input-head position $p<n+2$, every head
encoding $h<(2s+1)^{k}$, every machine case, and every wire index $v$ less than the
total width, where $k$ is the number of work tapes of $M$.
\end{definition}

\begin{definition}[Configuration adjacency formula]
\lean{Complexity.ConfigCodec.adjacency}
\label{Complexity.ConfigCodec.adjacency}
\leanok
\uses{Complexity.ConfigCodec.inputEquality, Complexity.ConfigCodec.transition}
For a finite Turing machine $M$ and natural numbers $n,s$, define the adjacency CNF
formula to be the concatenation of the input-equality formula and the transition formula
for these parameters.
\end{definition}

\begin{definition}[Turing machine configuration row]
\lean{Complexity.CodecRow}
\label{Complexity.CodecRow}
\leanok
\uses{Turing.FinTM, Turing.OutSummary}
For a finite-state Turing machine $M$ over the Boolean alphabet and natural numbers $n$
and $s$, a configuration row consists of an $n$-bit input, an optional state of $M$, an
input-head position among $n+2$ positions, a Boolean-or-blank symbol for each work tape
at each of $2s+1$ cells, a head position among those cells for each work tape, and an
output summary.
\end{definition}

\begin{definition}[Configuration encoding bit positions]
\lean{Complexity.CodecBit}
\label{Complexity.CodecBit}
\leanok
\uses{Turing.FinTM}
For a finite-state Turing machine $M$ and natural numbers $n$ and $s$, the bit positions
are tagged as one of the following: an input coordinate among $n$ coordinates; an
input-head coordinate among $n+2$ coordinates; a state choice, including the halting
choice; a work-tape coordinate consisting of a tape, one of $2s+1$ cells, and one of
three slots; or one of two output-summary coordinates.
\end{definition}

\begin{theorem}[Number of codec bits]
\lean{Complexity.codecBit_card}
\label{Complexity.codecBit_card}
\leanok
\uses{Complexity.CodecBit, Turing.FinTM}
\difficulty{3}
Let $M$ be a finite-state Turing machine over the Boolean alphabet, with $k$ work tapes
and state set $Q$, and let $n,s\ge 0$. The number of codec-bit positions is
$n+(n+2)+(|Q|+1)+3k(2s+1)+2$.
\end{theorem}

\begin{definition}[Configuration encoding coefficient]
\lean{Complexity.codecCoefficient}
\label{Complexity.codecCoefficient}
\leanok
\uses{Turing.FinTM}
For a finite-state Turing machine $M$, define its encoding coefficient to be the number
of states of $M$, plus six times its number of work tapes, plus five.
\end{definition}

\begin{theorem}[Codec-bit count bound]
\lean{Complexity.codecBit_card_le}
\label{Complexity.codecBit_card_le}
\leanok
\uses{Complexity.CodecBit, Complexity.codecBit_card, Complexity.codecCoefficient, Turing.FinTM}
\difficulty{3}
Let $M$ be a finite-state Turing machine over the Boolean alphabet, with $k$ work tapes
and state set $Q$, and let $n,s,C$ be nonnegative integers. If $|Q|+6k+5\le C$, then the
number of codec-bit positions is at most $C(s+n+1)$.
\end{theorem}

\begin{definition}[Bits of an encoding row]
\lean{Complexity.codecRowBits}
\label{Complexity.codecRowBits}
\leanok
\uses{Complexity.CodecBit, Complexity.CodecRow, Turing.FinTM}
Given a configuration encoding row, assign a Boolean value to each tagged bit position
as follows. At an input coordinate use the corresponding input bit; at an input-head or
state coordinate use $1$ exactly when that coordinate is the row's selected position or
state. For a work-tape coordinate, the first and second slots indicate, respectively,
that the cell contains false or true, and the third indicates that the work-tape head is
at that cell. At an output-summary coordinate, the bit is $1$ exactly when the row's
summary is dead for the first coordinate or accepting for the second.
\end{definition}

\begin{theorem}[Boolean symbols determined by their indicators]
\lean{Complexity.codec_symbol_ext}
\label{Complexity.codec_symbol_ext}
\leanok
\difficulty{4}
If $u,v \in \{\text{none},\text{false},\text{true}\}$ have the same truth value for the
predicates $u=\text{some false}$ and $v=\text{some false}$, and also the same truth
value for $u=\text{some true}$ and $v=\text{some true}$, then $u=v$.
\end{theorem}

\begin{theorem}[Injectivity of the row bit encoding]
\lean{Complexity.codecRowBits_injective}
\label{Complexity.codecRowBits_injective}
\leanok
\uses{Complexity.codecRowBits, Complexity.codec_symbol_ext, Turing.FinTM}
\difficulty{5}
For any finite Turing machine, input length, and window size, the map that sends a codec
row to its Boolean values on all codec-bit coordinates is injective.
\end{theorem}

\begin{definition}[Encoding row of a configuration]
\lean{Complexity.codecRowOf}
\label{Complexity.codecRowOf}
\leanok
\uses{Complexity.CodecRow, Turing.Cfg, Turing.FinTM, Turing.MultiTapeTM.ConfigCount.coreCode, Turing.outSummary}
Let $M$ be a finite-state Turing machine, let $n,s$ be natural numbers, let $x$ be a
Boolean string of length $n$, and let $c$ be a configuration of $M$ on input $x$. The
encoding row of $c$ records the bounded configuration code of $c$ at parameter $s$,
together with the input bits of $x$ (using false for any out-of-range lookup), and the
output summary determined by the output string of $c$.
\end{definition}

\begin{theorem}[Recovery of the input from a row]
\lean{Complexity.codecRowOf_input_eq}
\label{Complexity.codecRowOf_input_eq}
\leanok
\uses{Complexity.codecRowOf, Turing.Cfg, Turing.FinTM}
\difficulty{4}
Let $M$ be a finite Turing machine, let $n,s$ be natural numbers, and let $x,y$ be
Boolean lists each of length $n$. For any configurations $c,d$ of $M$ on those
respective inputs, if the input fields of their codec rows are equal, then $x=y$.
\end{theorem}

\begin{theorem}[Equality of rows and configuration cores]
\lean{Complexity.codecRowOf_eq_iff}
\label{Complexity.codecRowOf_eq_iff}
\leanok
\uses{Complexity.codecRowOf, Turing.Cfg, Turing.Cfg.InWindow, Turing.Cfg.InWindow.core_bounds, Turing.FinTM, Turing.MultiTapeTM.ConfigCount.core, Turing.MultiTapeTM.ConfigCount.coreCode, Turing.MultiTapeTM.ConfigCount.coreCode_inj, Turing.NDTM.coreSum, Turing.outSummary}
\difficulty{5}
Let $M$ be a finite Turing machine, let $n,s$ be natural numbers, and let $x$ be a
Boolean list of length $n$. For configurations $c,d$ of $M$ on input $x$ that are both
in the window of radius $s$, their codec rows are equal if and only if their cores
together with their output summaries are equal.
\end{theorem}

\begin{definition}[Boolean word from a finite index set]
\lean{Complexity.codecWord}
\label{Complexity.codecWord}
\leanok
For a finite type $\alpha$, a natural number $L$, and a Boolean-valued function $a$ on
$\alpha$, form a Boolean word of length $L$. Its coordinate $i$ is the value of $a$ at
the element with index $i$ in the standard enumeration of $\alpha$ when $i<|\alpha|$,
and is false otherwise.
\end{definition}

\begin{theorem}[Length of the Boolean code]
\lean{Complexity.codecWord_length}
\label{Complexity.codecWord_length}
\leanok
\uses{Complexity.codecWord}
\difficulty{1}
Let $\alpha$ be a finite type, let $L$ be a natural number, and let
$a\colon\alpha\to\{\mathrm{false},\mathrm{true}\}$. The Boolean list encoding $a$ with
length parameter $L$ has length exactly $L$.
\end{theorem}

\begin{theorem}[Recovering a Boolean value from its code]
\lean{Complexity.codecWord_get}
\label{Complexity.codecWord_get}
\leanok
\uses{Complexity.codecWord}
\difficulty{3}
Let $\alpha$ be a finite type, let $L$ be a natural number at least $|\alpha|$, and let
$a\colon\alpha\to\{\mathrm{false},\mathrm{true}\}$. Encode $a$ as a Boolean list of
length $L$ by placing its values in the order given by an enumeration of $\alpha$ and
padding with false values. Then the list entry at the position of any $i\in\alpha$ is
exactly $a(i)$.
\end{theorem}

\begin{theorem}[Injectivity of the Boolean Word Encoding]
\lean{Complexity.codecWord_eq_iff}
\label{Complexity.codecWord_eq_iff}
\leanok
\uses{Complexity.codecWord, Complexity.codecWord_get}
\difficulty{4}
Let $\alpha$ be a finite type, let $L\in\mathbb{N}$ be at least the cardinality of
$\alpha$, and let $a,b:\alpha\to\{\mathrm{false},\mathrm{true}\}$. Their length-$L$
encoded Boolean words are equal exactly when $a=b$.
\end{theorem}

\begin{definition}[Configuration Boolean encoding]
\lean{Complexity.configurationCodec}
\label{Complexity.configurationCodec}
\leanok
\uses{Complexity.codecRowBits, Complexity.codecRowOf, Complexity.codecWord, Turing.Cfg, Turing.FinTM}
Let $M$ be a finite-state Turing machine, let $C,n,s$ be natural numbers, let $x$ be a
Boolean string, and let $c$ be a configuration of $M$ on input $x$. If $x$ has length
$n$, encode the bits of its configuration row as a Boolean word of length $C(s+n+1)$. If
$x$ does not have length $n$, use the all-false Boolean word of that length.
\end{definition}

\begin{theorem}[Configuration code length]
\lean{Complexity.configurationCodec_length}
\label{Complexity.configurationCodec_length}
\leanok
\uses{Complexity.codecWord_length, Complexity.configurationCodec, Turing.Cfg, Turing.FinTM}
\difficulty{4}
For a finite-state Turing machine with Boolean input alphabet, any natural numbers
$C,n,s$, any Boolean input string $x$, and any configuration on that input, the
configuration code has length exactly $C(s+n+1)$.
\end{theorem}

\begin{theorem}[Input Recovery from Configuration Encodings]
\lean{Complexity.configurationCodec_input_eq}
\label{Complexity.configurationCodec_input_eq}
\leanok
\uses{Complexity.codecBit_card_le, Complexity.codecCoefficient, Complexity.codecRowBits_injective, Complexity.codecRowOf_input_eq, Complexity.codecWord_eq_iff, Complexity.configurationCodec, Turing.Cfg, Turing.FinTM}
\difficulty{3}
Let $M$ be a finite-state Turing machine with $k$ work tapes, and let $n,s,C\in\mathbb
N$ with $C\ge |Q|+6k+5$, where $Q$ is its state set. If $x$ and $y$ are Boolean inputs
of length $n$, and configurations $c$ on $x$ and $d$ on $y$ have equal configuration
encodings, then $x=y$.
\end{theorem}

\begin{theorem}[Injectivity of the Configuration Encoding]
\lean{Complexity.configurationCodec_eq_iff}
\label{Complexity.configurationCodec_eq_iff}
\leanok
\uses{Complexity.codecBit_card_le, Complexity.codecCoefficient, Complexity.codecRowBits_injective, Complexity.codecRowOf_eq_iff, Complexity.codecWord_eq_iff, Complexity.configurationCodec, Turing.Cfg, Turing.Cfg.InWindow, Turing.FinTM, Turing.NDTM.coreSum}
\difficulty{3}
Let $M$ be a finite-state Turing machine with $k$ work tapes, and let $n,s,C\in\mathbb
N$ with $C\ge |Q|+6k+5$, where $Q$ is its state set. For any Boolean input $x$ of length
$n$ and configurations $c,d$ on $x$ whose work-tape heads lie within distance $s$ of the
origin and whose work tapes are blank outside that interval, their configuration
encodings are equal exactly when their states, input-head positions, work-tape contents,
work-tape head positions, and output summaries are all equal.
\end{theorem}

\begin{theorem}[Evaluation of a flattened CNF list]
\lean{Complexity.codecCnf_eval_flatMap}
\label{Complexity.codecCnf_eval_flatMap}
\leanok
\difficulty{1}
Let $X$ be a list, let each element be assigned a CNF over a fixed variable type, and
let $a$ assign a Boolean value to each variable. The concatenation of the assigned CNFs
evaluates to true under $a$ if and only if every CNF assigned to an element occurring in
$X$ evaluates to true under $a$.
\end{theorem}

\begin{theorem}[Evaluation of a mapped CNF list]
\lean{Complexity.codecCnf_eval_map}
\label{Complexity.codecCnf_eval_map}
\leanok
\difficulty{1}
Let $X$ be a list of objects, let each object $i$ determine a CNF clause $f(i)$ over
variables in a set $A$, and let $a:A\to\{\mathrm{false},\mathrm{true}\}$ be an
assignment. The conjunction of the clauses in the list obtained by applying $f$ to every
element of $X$ evaluates to true if and only if every clause $f(i)$ for $i\in X$
evaluates to true.
\end{theorem}

\begin{definition}[One-hot Boolean clauses]
\lean{Complexity.codecOneHot}
\label{Complexity.codecOneHot}
\leanok
For a finite type $\alpha$ with decidable equality, define a conjunctive normal form
over variables indexed by $\alpha$ whose first clause contains every variable
positively, and which also contains, for each pair of distinct variables, a clause
containing both variables negatively.
\end{definition}

\begin{theorem}[Evaluation of the one-hot encoding]
\lean{Complexity.codecOneHot_eval}
\label{Complexity.codecOneHot_eval}
\leanok
\uses{Complexity.codecOneHot}
\difficulty{5}
For every finite type $\alpha$ and Boolean assignment $a$ on $\alpha$, the one-hot CNF
encoding evaluates to true if and only if there is some $j\in\alpha$ such that $a(i)$ is
true exactly when $i=j$, for every $i\in\alpha$.
\end{theorem}

\begin{definition}[Validity of a machine configuration encoding]
\lean{Complexity.codecBitsValid}
\label{Complexity.codecBitsValid}
\leanok
\uses{Complexity.CodecBit, Turing.FinTM}
Let $M$ be a finite Turing machine over the Boolean alphabet, let $n,s\in\mathbb{N}$,
and assign a Boolean value $a$ to each bit of the encoding, whose bits represent input
positions, an input-head position, a state, work-tape cells, and a summary. The
assignment is valid when exactly one input-head position and exactly one state are
selected; on each work tape exactly one of the $2s+1$ positions is selected; at each
work-tape position the two bits indexed by $0$ and $1$ are not both true; and the two
summary bits are not both true.
\end{definition}

\begin{theorem}[Validity of encoded row bits]
\lean{Complexity.codecRowBits_valid}
\label{Complexity.codecRowBits_valid}
\leanok
\uses{Complexity.CodecRow, Complexity.codecBitsValid, Complexity.codecRowBits, Turing.FinTM}
\difficulty{4}
For any finite Turing machine $M$, natural numbers $n,s$, and codec row $r$, its Boolean
encoding has exactly one true input-head coordinate and exactly one true state
coordinate. For each work tape, exactly one head-position coordinate is true; at each
work-tape cell the two symbol coordinates are not both true; and the two output-summary
coordinates are not both true.
\end{theorem}

\begin{theorem}[Characterization of valid codec bits]
\lean{Complexity.codecBitsValid_iff}
\label{Complexity.codecBitsValid_iff}
\leanok
\uses{Complexity.CodecBit, Complexity.CodecRow, Complexity.codecBitsValid, Complexity.codecRowBits, Complexity.codecRowBits_valid, Turing.FinTM}
\difficulty{6}
Let $M$ be a finite-state Turing machine over the Boolean alphabet, let $n,s\ge0$, and
assign a Boolean value to every codec-bit position. This assignment satisfies the
validity conditions exactly when it is the bit encoding of some row consisting of an
input word of length $n$, an optional state, an input-head position among $n+2$
positions, work-tape symbols and head positions on each of $k$ tapes across $2s+1$
cells, and an output summary. The validity conditions require exactly one input-head
position and state, exactly one head position per work tape, never simultaneously
marking a cell as containing both Boolean symbols, and never marking both output-summary
bits true.
\end{theorem}

\begin{definition}[CNF encoding of configuration validity]
\lean{Complexity.codecValidity}
\label{Complexity.codecValidity}
\leanok
\uses{Complexity.CodecBit, Complexity.codecOneHot, Turing.FinTM}
Let $M$ be a finite Turing machine over the Boolean alphabet, and let
$n,s\in\mathbb{N}$. The configuration-validity CNF over the encoding bits consists of
one-hot constraints selecting an input-head position and a state, one one-hot constraint
selecting a position on each work tape, clauses forbidding the two designated bits at
any work-tape position from both being true, and a clause forbidding both summary bits
from being true.
\end{definition}

\begin{theorem}[Validity Formula Characterization]
\lean{Complexity.codecValidity_eval}
\label{Complexity.codecValidity_eval}
\leanok
\uses{Complexity.CodecBit, Complexity.CodecRow, Complexity.codecBitsValid, Complexity.codecBitsValid_iff, Complexity.codecCnf_eval_flatMap, Complexity.codecCnf_eval_map, Complexity.codecOneHot_eval, Complexity.codecRowBits, Complexity.codecValidity, Turing.FinTM}
\difficulty{3}
Let $M$ be a finite-state Turing machine with $k$ work tapes, let $n,s\in\mathbb{N}$,
and assign a Boolean value to every encoded bit. The validity CNF evaluates to true
exactly when that assignment is the bit encoding of some codec row.
\end{theorem}

\begin{theorem}[Decoding a position code]
\lean{Complexity.codec_posCode}
\label{Complexity.codec_posCode}
\leanok
\uses{Turing.MultiTapeTM.ConfigCount.posCode}
\difficulty{3}
For every natural number $s$ and every $j\in\{0,1,\ldots,2s\}$, the position code with
bound $s$ of the integer $j-s$ is $j$.
\end{theorem}

\begin{theorem}[Realization of every codec row]
\lean{Complexity.codecRow_realize}
\label{Complexity.codecRow_realize}
\leanok
\uses{Complexity.CodecRow, Complexity.codecRowOf, Complexity.codec_posCode, Turing.Cfg, Turing.Cfg.InWindow, Turing.FinTM, Turing.MultiTapeTM.ConfigCount.coreCode, Turing.MultiTapeTM.ConfigCount.posCode, Turing.outSummary}
\difficulty{5}
For every finite Turing machine $M$, natural numbers $n,s$, and codec row $r$ for
$M,n,s$, there exist a Boolean list $x$ of length $n$ and a configuration $c$ of $M$ on
input $x$ such that $c$ is in the window of radius $s$ and its codec row is $r$.
\end{theorem}

\begin{theorem}[False padding in the Boolean code]
\lean{Complexity.codecWord_padding}
\label{Complexity.codecWord_padding}
\leanok
\uses{Complexity.codecWord}
\difficulty{4}
Let $\alpha$ be a finite type, let $L$ be a natural number, and let
$a\colon\alpha\to\{\mathrm{false},\mathrm{true}\}$. For every natural number $j\ge
|\alpha|$, the entry at position $j$ of the Boolean list encoding $a$ with length
parameter $L$, using false as the default value beyond the list, is false.
\end{theorem}

\begin{theorem}[Reconstructing a Boolean list from its code]
\lean{Complexity.codecWord_reconstruct}
\label{Complexity.codecWord_reconstruct}
\leanok
\uses{Complexity.codecWord, Complexity.codecWord_length}
\difficulty{4}
Let $\alpha$ be a finite type, let $L$ be a natural number, and let $w$ be a Boolean
list of length $L$. Suppose every position $j$ with $|\alpha|\le j<L$ contains false.
Then encoding the function that assigns to each $i\in\alpha$ the entry of $w$ at the
position of $i$ recovers $w$ exactly.
\end{theorem}

\begin{definition}[Padding clauses]
\lean{Complexity.codecPadding}
\label{Complexity.codecPadding}
\leanok
For natural numbers $first$ and $L$, form a CNF with one negative unit clause for each
index $j$ satisfying $first\le j<L$. Thus every such variable is forced to be false; if
$L\le first$, the CNF has no clauses.
\end{definition}

\begin{theorem}[Evaluation of the padding encoding]
\lean{Complexity.codecPadding_eval}
\label{Complexity.codecPadding_eval}
\leanok
\uses{Complexity.codecPadding}
\difficulty{4}
For all natural numbers $\mathrm{first},L$ and every Boolean assignment $a$ on the
natural numbers, the padding CNF encoding evaluates to true if and only if $a(j)$ is
false for every natural number $j$ with $\mathrm{first}\le j<L$.
\end{theorem}

\begin{definition}[Padded configuration-validity CNF]
\lean{Complexity.configurationValidity}
\label{Complexity.configurationValidity}
\leanok
\uses{Complexity.CodecBit, Complexity.codecPadding, Complexity.codecValidity, Turing.FinTM}
Let $M$ be a finite Turing machine over the Boolean alphabet, and let
$C,n,s\in\mathbb{N}$. Relabel the configuration-validity CNF by assigning its encoding
bits consecutive natural-number indices, then append negative unit clauses forcing every
index from the number of encoding bits up to, but not including, $C(s+n+1)$ to be false.
\end{definition}

\begin{theorem}[Characterization of valid configuration codes]
\lean{Complexity.configurationValidity_eval}
\label{Complexity.configurationValidity_eval}
\leanok
\uses{Complexity.codecBit_card_le, Complexity.codecCoefficient, Complexity.codecPadding_eval, Complexity.codecRowOf, Complexity.codecRow_realize, Complexity.codecValidity_eval, Complexity.codecWord_get, Complexity.codecWord_padding, Complexity.codecWord_reconstruct, Complexity.configurationCodec, Complexity.configurationValidity, Turing.Cfg, Turing.Cfg.InWindow, Turing.FinTM}
\difficulty{5}
Let $M$ be a finite-state Turing machine with Boolean input alphabet, and let $C,n,s$ be
natural numbers such that $C$ is at least the number of states plus $6$ times the number
of work tapes plus $5$. For every Boolean word $w$ of length $C(s+n+1)$, the
configuration-validity CNF formula evaluates to true when variable $j$ is assigned the
bit of $w$ at position $j$ (and false beyond the word) if and only if there are a
Boolean input string $x$ of length $n$ and a configuration $c$ on $x$, with all
work-tape heads within distance $s$ of the origin and all work-tape cells outside that
window blank, such that $w$ is the configuration code of $c$.
\end{theorem}

\begin{definition}[Acceptance constraints for the encoding]
\lean{Complexity.codecAcceptance}
\label{Complexity.codecAcceptance}
\leanok
\uses{Complexity.CodecBit, Turing.FinTM}
Let $M$ be a finite Turing machine over the Boolean alphabet, and let
$n,s\in\mathbb{N}$. The acceptance CNF consists of unit clauses requiring the encoded
state to be absent and the summary bit indexed by false to be true.
\end{definition}

\begin{theorem}[Acceptance formula evaluation]
\lean{Complexity.codecAcceptance_eval}
\label{Complexity.codecAcceptance_eval}
\leanok
\uses{Complexity.CodecRow, Complexity.codecAcceptance, Complexity.codecRowBits, Turing.FinTM}
\difficulty{1}
For every finite Turing machine $M$, natural numbers $n,s$, and encoded configuration
row $r$, the CNF formula consisting of the unit clauses requiring that the state is
halted and the output summary is accepting evaluates to true under the Boolean encoding
of $r$ if and only if $r$ is halted and its output summary is accepting.
\end{theorem}

\begin{definition}[Relabeled acceptance constraints]
\lean{Complexity.configurationAcceptance}
\label{Complexity.configurationAcceptance}
\leanok
\uses{Complexity.CodecBit, Complexity.codecAcceptance, Turing.FinTM}
Let $M$ be a finite Turing machine over the Boolean alphabet, and let
$n,s\in\mathbb{N}$. Relabel the two acceptance unit clauses by assigning the encoding
bits consecutive natural-number indices; the clauses require the state to be absent and
the summary bit indexed by false to be true.
\end{definition}

\begin{theorem}[Evaluation of the Acceptance Formula]
\lean{Complexity.configurationAcceptance_eval}
\label{Complexity.configurationAcceptance_eval}
\leanok
\uses{Complexity.codecAcceptance_eval, Complexity.codecBit_card_le, Complexity.codecCoefficient, Complexity.codecRowOf, Complexity.codecWord_get, Complexity.configurationAcceptance, Complexity.configurationCodec, Turing.Cfg, Turing.FinTM, Turing.outSummary}
\difficulty{3}
Let $M$ be a finite-state Turing machine with $k$ work tapes, and let $n,s,C\in\mathbb
N$ with $C\ge |Q|+6k+5$, where $Q$ is its state set. For any Boolean input $x$ of length
$n$ and any configuration $c$ on that input, the acceptance formula evaluated on the
configuration's encoded bit string is true exactly when $c$ is halted and its output is
exactly the one-symbol string $[\mathrm{true}]$.
\end{theorem}

\begin{definition}[CNF encoding cost]
\lean{Complexity.codecCnfCost}
\label{Complexity.codecCnfCost}
\leanok
For a CNF formula $\varphi$, define its cost to be the sum, over all clauses of
$\varphi$, of one plus the number of literals in the clause.
\end{definition}

\begin{theorem}[Invariance of CNF cost under relabeling]
\lean{Complexity.codecCnfCost_relabel}
\label{Complexity.codecCnfCost_relabel}
\leanok
\uses{Complexity.codecCnfCost}
\difficulty{1}
For any CNF and any function from its variable type to another type, relabeling the
variables by that function leaves the CNF cost unchanged, where cost is the sum over
clauses of one plus the clause length.
\end{theorem}

\begin{theorem}[Cost of concatenated CNFs]
\lean{Complexity.codecCnfCost_append}
\label{Complexity.codecCnfCost_append}
\leanok
\uses{Complexity.codecCnfCost}
\difficulty{1}
For any two CNFs over the same variable type, the cost of their concatenation equals the
sum of their costs, where each CNF's cost is the sum over its clauses of one plus the
clause length.
\end{theorem}

\begin{theorem}[Cost bound for a flattened CNF list]
\lean{Complexity.codecCnfCost_flatMap_le}
\label{Complexity.codecCnfCost_flatMap_le}
\leanok
\uses{Complexity.codecCnfCost, Complexity.codecCnfCost_append}
\difficulty{4}
Let $X$ be a list and assign to each element $x$ a CNF over a fixed variable type. If
every CNF assigned to an element occurring in $X$ has cost at most $B$, then the cost of
the CNF obtained by concatenating these CNFs in list order is at most $|X|B$.
\end{theorem}

\begin{theorem}[One-hot encoding cost bound]
\lean{Complexity.codecOneHot_cost}
\label{Complexity.codecOneHot_cost}
\leanok
\uses{Complexity.codecCnfCost, Complexity.codecCnfCost_append, Complexity.codecCnfCost_flatMap_le, Complexity.codecOneHot}
\difficulty{5}
For every finite type $\alpha$, the cost of the one-hot CNF encoding on $\alpha$ is at
most $|\alpha|+1+3|\alpha|^2$, where $|\alpha|$ is the number of elements of $\alpha$.
\end{theorem}

\begin{theorem}[Variable bound after relabeling]
\lean{Complexity.codecCnf_numVars}
\label{Complexity.codecCnf_numVars}
\leanok
\uses{Std.Sat.CNF.numVars}
\difficulty{4}
Let $\varphi$ be a CNF formula over a type of variables, let $w$ map each variable to a
natural-number index, and let $N$ be a natural number. Suppose every literal variable
occurring in every clause of $\varphi$ is mapped by $w$ to an index strictly less than
$N$. Then the variable count of the formula obtained by relabeling $\varphi$ with $w$ is
at most $N$.
\end{theorem}

\begin{theorem}[Length bound for CNF serialization]
\lean{Complexity.codecCnf_serialize}
\label{Complexity.codecCnf_serialize}
\leanok
\uses{Complexity.codecCnfCost, Std.Sat.CNF.serialize, Std.Sat.CNF.serializeClause, Std.Sat.CNF.serializeLit}
\difficulty{5}
Let $\varphi$ be a CNF formula over a type of variables, let $w$ assign a natural-number
index to each variable, and let $N$ be a natural number. If every variable in every
clause of $\varphi$ is mapped to an index less than $N$, then the length of the binary
serialization of the relabeled formula is at most $\operatorname{cost}(\varphi)(N+3)+1$,
where $\operatorname{cost}(\varphi)$ is the sum, over its clauses $C$, of $|C|+1$.
\end{theorem}

\begin{definition}[Validity coefficient]
\lean{Complexity.codecValidityCoefficient}
\label{Complexity.codecValidityCoefficient}
\leanok
\uses{Turing.FinTM}
For a finite Turing machine $M$ with $k$ work tapes and $q$ states, where $q$ denotes
the number of states plus one, define its validity coefficient to be $3q^2+q+9+21k$.
\end{definition}

\begin{theorem}[Raw Validity Formula Cost Bound]
\lean{Complexity.codecValidity_cost_raw}
\label{Complexity.codecValidity_cost_raw}
\leanok
\uses{Complexity.codecCnfCost, Complexity.codecCnfCost_append, Complexity.codecCnfCost_flatMap_le, Complexity.codecCnfCost_relabel, Complexity.codecOneHot, Complexity.codecOneHot_cost, Complexity.codecValidity, Turing.FinTM}
\difficulty{5}
Let $M$ be a finite-state Turing machine with $k$ work tapes, and let
$n,s\in\mathbb{N}$. The validity CNF cost is at most
$((n+2)+1+3(n+2)^2)+((|\mathrm{State}(M)|+1)+1+3(|\mathrm{State}(M)|+1)^2)+k((2s+1)+1+3(2s+1)^2)+3k(2s+1)+3$.
\end{theorem}

\begin{theorem}[Validity Formula Cost]
\lean{Complexity.codecValidity_cost}
\label{Complexity.codecValidity_cost}
\leanok
\uses{Complexity.codecCnfCost, Complexity.codecValidity, Complexity.codecValidityCoefficient, Complexity.codecValidity_cost_raw, Turing.FinTM}
\difficulty{5}
Let $M$ be a finite-state Turing machine with $k$ work tapes, and let
$n,s\in\mathbb{N}$. The cost of the validity CNF is at most
$\bigl(3(|\mathrm{State}(M)|+1)^2+(|\mathrm{State}(M)|+1)+9+21k\bigr)(s+n+2)^2$.
\end{theorem}

\begin{theorem}[Padding encoding cost bound]
\lean{Complexity.codecPadding_cost}
\label{Complexity.codecPadding_cost}
\leanok
\uses{Complexity.codecCnfCost, Complexity.codecPadding}
\difficulty{3}
For all natural numbers $\mathrm{first}$ and $L$, the cost of the padding CNF encoding
is at most $2L$.
\end{theorem}

\begin{theorem}[Variable indices in the validity formula]
\lean{Complexity.configurationValidity_literals}
\label{Complexity.configurationValidity_literals}
\leanok
\uses{Complexity.CodecBit, Complexity.codecBit_card_le, Complexity.codecCoefficient, Complexity.codecPadding, Complexity.configurationValidity, Turing.FinTM}
\difficulty{4}
Let $M$ be a finite-state Turing machine with Boolean input alphabet, and let $C,n,s$ be
natural numbers such that $C$ is at least the number of states plus $6$ times the number
of work tapes plus $5$. Every variable appearing in every clause of the
configuration-validity CNF formula has index less than $C(s+n+1)$.
\end{theorem}

\begin{theorem}[Variable count of the validity formula]
\lean{Complexity.configurationValidity_numVars}
\label{Complexity.configurationValidity_numVars}
\leanok
\uses{Complexity.codecCnf_numVars, Complexity.codecCoefficient, Complexity.configurationValidity, Complexity.configurationValidity_literals, Std.Sat.CNF.numVars, Turing.FinTM}
\difficulty{1}
Let $M$ be a finite-state Turing machine with Boolean input alphabet, and let $C,n,s$ be
natural numbers such that $C$ is at least the number of states plus $6$ times the number
of work tapes plus $5$. The configuration-validity CNF formula has at most $C(s+n+1)$
variables.
\end{theorem}

\begin{theorem}[Raw Configuration Encoding Length]
\lean{Complexity.configurationValidity_size_raw}
\label{Complexity.configurationValidity_size_raw}
\leanok
\uses{Complexity.CodecBit, Complexity.codecCnfCost, Complexity.codecCnfCost_append, Complexity.codecCnfCost_relabel, Complexity.codecCnf_serialize, Complexity.codecCoefficient, Complexity.codecPadding_cost, Complexity.codecValidityCoefficient, Complexity.codecValidity_cost, Complexity.configurationValidity, Complexity.configurationValidity_literals, Std.Sat.CNF.serialize, Turing.FinTM}
\difficulty{4}
Let $M$ be a finite Turing machine with $k$ work tapes and state set $Q$, and let
$C,n,s$ be natural numbers. Write $A=|Q|+6k+5$ and $B=3(|Q|+1)^2+(|Q|+1)+9+21k$; assume
$A\le C$. The serialized configuration-validity CNF for parameters $C,n,s$ has length at
most $\bigl(B(s+n+2)^2+2C(s+n+1)\bigr)\bigl(C(s+n+1)+3\bigr)+1$.
\end{theorem}

\begin{theorem}[Configuration Acceptance Bounds]
\lean{Complexity.configurationAcceptance_bounds}
\label{Complexity.configurationAcceptance_bounds}
\leanok
\uses{Complexity.CodecBit, Complexity.codecAcceptance, Complexity.codecBit_card_le, Complexity.codecCnfCost, Complexity.codecCnf_numVars, Complexity.codecCnf_serialize, Complexity.codecCoefficient, Complexity.configurationAcceptance, Std.Sat.CNF.numVars, Std.Sat.CNF.serialize, Turing.FinTM}
\difficulty{3}
Let $M$ be a finite-state Turing machine with $k$ work tapes, and let $n,s,C\in\mathbb
N$. If $C\ge |Q|+6k+5$, where $Q$ is the state set of $M$, then the variable count of
the CNF formula encoding acceptance is at most $C(s+n+1)$, and the length of its binary
serialization is at most $4\bigl(C(s+n+1)+3\bigr)+1$.
\end{theorem}

\begin{theorem}[A polynomial size bound]
\lean{Complexity.codec_size_normalize}
\label{Complexity.codec_size_normalize}
\leanok
\difficulty{5}
For natural numbers $A,C,t$ with $C\ge1$ and $t\ge1$, the following inequality holds:
$\bigl(A(t+1)^2+2Ct\bigr)(Ct+3)+1\le(3A+7)C^2(t+1)^3$.
\end{theorem}

\begin{theorem}[Configuration Validity Encoding Length]
\lean{Complexity.configurationValidity_size}
\label{Complexity.configurationValidity_size}
\leanok
\uses{Complexity.codecCoefficient, Complexity.codecValidityCoefficient, Complexity.codec_size_normalize, Complexity.configurationValidity, Complexity.configurationValidity_size_raw, Std.Sat.CNF.serialize, Turing.FinTM}
\difficulty{3}
Let $M$ be a finite Turing machine with $k$ work tapes and state set $Q$, and let
$C,n,s$ be natural numbers. Write $A=|Q|+6k+5$ and $B=3(|Q|+1)^2+(|Q|+1)+9+21k$; assume
$A\le C$. The serialized configuration-validity CNF for parameters $C,n,s$ has length at
most $(3B+7)C^2(s+n+2)^3$.
\end{theorem}

\begin{theorem}[Size of the Acceptance Formula]
\lean{Complexity.configurationAcceptance_size}
\label{Complexity.configurationAcceptance_size}
\leanok
\uses{Complexity.codecCoefficient, Complexity.configurationAcceptance, Complexity.configurationAcceptance_bounds, Std.Sat.CNF.serialize, Turing.FinTM}
\difficulty{4}
Let $M$ be a finite-state Turing machine with $k$ work tapes, and let $n,s,C\in\mathbb
N$ with $C\ge |Q|+6k+5$, where $Q$ is its state set. The length of the binary
serialization of the acceptance formula is at most $13C^2(s+n+2)^3$.
\end{theorem}

\begin{theorem}[CNF representation of a local Boolean function]
\lean{Complexity.codecLocalCheck}
\label{Complexity.codecLocalCheck}
\leanok
\uses{Complexity.codecCnfCost, Complexity.codecCnfCost_flatMap_le, Complexity.codecCnfCost_relabel, Complexity.exists_cnf_boolFun}
\difficulty{4}
Let $A$ be an inhabited type, let $m$ be a natural number, let $w$ assign an element of
$A$ to each of the $m$ input positions, and let $f$ be a Boolean function of $m$ Boolean
inputs. There is a CNF formula over variables in $A$ with cost at most $2^m(m+1)$ such
that, for every Boolean assignment to $A$, the formula evaluates to $f$ applied to the
assignment's values at the positions selected by $w$.
\end{theorem}

\begin{definition}[Configuration codec positions]
\lean{Complexity.CodecPositions}
\label{Complexity.CodecPositions}
\leanok
\uses{Turing.FinTM}
For a finite Turing machine $M$ with Boolean alphabet and natural numbers $n$ and $s$, a
codec position is a pair consisting of an input-head position among $n+2$ positions and,
for each of the $M$'s work tapes, a position among $2s+1$ positions.
\end{definition}

\begin{theorem}[Bound on the number of codec positions]
\lean{Complexity.codecPositions_card_le}
\label{Complexity.codecPositions_card_le}
\leanok
\uses{Complexity.CodecPositions, Turing.FinTM}
\difficulty{3}
Let $M$ be a finite-state Turing machine with Boolean symbols, and let $n,s$ be natural
numbers. The number of codec positions is at most $2^k(s+n+2)^{k+1}$, where $k$ is the
number of work tapes of $M$.
\end{theorem}

\begin{definition}[Selected codec row bit]
\lean{Complexity.codecPairBits}
\label{Complexity.codecPairBits}
\leanok
\uses{Complexity.CodecBit, Complexity.CodecRow, Complexity.codecRowBits, Turing.FinTM}
Let $M$ be a finite Turing machine with Boolean alphabet, let $n,s$ be natural numbers,
let $r,t$ be codec rows, and let $p$ consist of a Boolean and a codec-bit coordinate.
The selected bit is the bit of $t$ at that coordinate when the Boolean is true, and the
bit of $r$ there when it is false.
\end{definition}

\begin{definition}[Wire number for a codec bit]
\lean{Complexity.codecPairWire}
\label{Complexity.codecPairWire}
\leanok
\uses{Complexity.CodecBit, Turing.FinTM}
Let $M$ be a finite Turing machine with Boolean alphabet, and let $C,n,s$ be natural
numbers. A Boolean paired with a codec-bit coordinate is assigned its finite-enumeration
index, with the offset $C(s+n+1)$ added when the Boolean is true and zero added when it
is false.
\end{definition}

\begin{theorem}[Pair-wire index bound]
\lean{Complexity.codecPairWire_lt}
\label{Complexity.codecPairWire_lt}
\leanok
\uses{Complexity.CodecBit, Complexity.codecBit_card_le, Complexity.codecCoefficient, Complexity.codecPairWire, Turing.FinTM}
\difficulty{4}
Let $M$ be a finite-state Turing machine with Boolean symbols, and let $C,n,s$ be
natural numbers. Define $D$ to be the number of states of $M$ plus $6$ times its number
of work tapes plus $5$. If $D\le C$, then every pair consisting of a Boolean value and a
codec bit has wire index less than $2C(s+n+1)$.
\end{theorem}

\begin{theorem}[Paired configuration code lookup]
\lean{Complexity.configurationCodec_pair_get}
\label{Complexity.configurationCodec_pair_get}
\leanok
\uses{Complexity.CodecBit, Complexity.codecBit_card_le, Complexity.codecCoefficient, Complexity.codecPairBits, Complexity.codecPairWire, Complexity.codecRowOf, Complexity.codecWord_get, Complexity.configurationCodec, Complexity.configurationCodec_length, Turing.Cfg, Turing.FinTM}
\difficulty{5}
Let $M$ be a finite-state Turing machine with Boolean input alphabet, and let $C,n,s$ be
natural numbers such that $C$ is at least the number of states plus $6$ times the number
of work tapes plus $5$. If Boolean strings $x,y$ each have length $n$, and $c,d$ are
configurations on those respective inputs, then for any choice of one of the two strings
and any coordinate of a configuration code, the bit at the corresponding position in the
concatenation of their configuration codes equals the bit specified by that chosen
configuration's code row at that coordinate.
\end{theorem}

\begin{definition}[Input-bit equality clauses]
\lean{Complexity.codecInputEquality}
\label{Complexity.codecInputEquality}
\leanok
\uses{Complexity.CodecBit, Turing.FinTM}
Let $M$ be a finite Turing machine with Boolean alphabet, and let $n,s$ be natural
numbers. For each input position $i<n$, form the two clauses expressing that the bit for
$i$ in the first row equals the bit for $i$ in the second row; the resulting conjunctive
normal form is the list of these clause pairs in increasing input-position order.
\end{definition}

\begin{theorem}[Input equality formula evaluation]
\lean{Complexity.codecInputEquality_eval}
\label{Complexity.codecInputEquality_eval}
\leanok
\uses{Complexity.CodecRow, Complexity.codecCnf_eval_flatMap, Complexity.codecInputEquality, Complexity.codecPairBits, Complexity.codecRowBits, Turing.FinTM}
\difficulty{4}
For every finite Turing machine $M$ over Booleans, natural numbers $n,s$, and encoded
rows $r,t$, the input-equality CNF formula evaluates to true on the assignment that
reads each literal from $r$ or $t$ according to its row selector if and only if the
input fields of $r$ and $t$ are equal.
\end{theorem}

\begin{theorem}[Input equality formula cost]
\lean{Complexity.codecInputEquality_cost}
\label{Complexity.codecInputEquality_cost}
\leanok
\uses{Complexity.codecCnfCost, Complexity.codecCnfCost_flatMap_le, Complexity.codecInputEquality, Turing.FinTM}
\difficulty{3}
For every finite Turing machine $M$ over Booleans and natural numbers $n,s$, the CNF
formula expressing equality of the two encoded input rows has cost at most $6n$.
\end{theorem}

\begin{definition}[Relabelled input equality formula]
\lean{Complexity.configurationInputEquality}
\label{Complexity.configurationInputEquality}
\leanok
\uses{Complexity.codecInputEquality, Complexity.codecPairWire, Turing.FinTM}
Let $M$ be a finite Turing machine with Boolean alphabet, and let $C,n,s$ be natural
numbers. Relabel every variable in the input-bit equality formula by the natural-number
wire assigned to its Boolean-and-codec-bit pair using offset parameter $C$.
\end{definition}

\begin{theorem}[Input inequality rejects paired codes]
\lean{Complexity.configurationInputEquality_reject}
\label{Complexity.configurationInputEquality_reject}
\leanok
\uses{Complexity.codecCoefficient, Complexity.codecInputEquality_eval, Complexity.codecRowOf_input_eq, Complexity.configurationCodec, Complexity.configurationCodec_pair_get, Complexity.configurationInputEquality, Turing.Cfg, Turing.FinTM}
\difficulty{4}
Let $M$ be a finite-state Turing machine with Boolean input alphabet, and let $C,n,s$ be
natural numbers such that $C$ is at least the number of states plus $6$ times the number
of work tapes plus $5$. If distinct Boolean strings $x,y$ both have length $n$, and
$c,d$ are configurations on these inputs, then the CNF formula enforcing equality of the
two inputs evaluates to false on the assignment given by the bits of the concatenated
configuration codes, with any out-of-range variable assigned false.
\end{theorem}

\begin{theorem}[Input equality formula bounds]
\lean{Complexity.configurationInputEquality_bounds}
\label{Complexity.configurationInputEquality_bounds}
\leanok
\uses{Complexity.codecCnf_numVars, Complexity.codecCnf_serialize, Complexity.codecCoefficient, Complexity.codecInputEquality, Complexity.codecInputEquality_cost, Complexity.codecPairWire, Complexity.codecPairWire_lt, Complexity.configurationInputEquality, Std.Sat.CNF.numVars, Std.Sat.CNF.serialize, Turing.FinTM}
\difficulty{4}
Let $M$ be a finite-state Turing machine with Boolean input alphabet, and let $C,n,s$ be
natural numbers such that $C$ is at least the number of states plus $6$ times the number
of work tapes plus $5$. The CNF formula enforcing equality of the input bits in a pair
of configuration codes has at most $2C(s+n+1)$ variables, and its binary serialization
has length at most $6n(2C(s+n+1)+3)+1$.
\end{theorem}

\begin{theorem}[Step determined by configuration code]
\lean{Complexity.configurationCodec_step_congr}
\label{Complexity.configurationCodec_step_congr}
\leanok
\uses{Complexity.codecCoefficient, Complexity.configurationCodec, Complexity.configurationCodec_eq_iff, Turing.Cfg, Turing.Cfg.InWindow, Turing.FinTM, Turing.MultiTapeTM.step, Turing.MultiTapeTM.toNDTM, Turing.NDTM.coreSum, Turing.NDTM.coreSum_stepWith}
\difficulty{3}
Let $M$ be a finite-state Turing machine with Boolean input alphabet, and let $C,n,s$ be
natural numbers such that $C$ is at least the number of states plus $6$ times the number
of work tapes plus $5$. Let $x$ be a Boolean string of length $n$, and let $c,d$ be
configurations on $x$ whose work-tape head positions lie within distance $s$ of the
origin and whose tape contents are blank outside that window. If $c$ and $d$ have the
same configuration code, then after one step of $M$ their states, input-head positions,
work-tape contents, work-tape head positions, and output summaries agree.
\end{theorem}

\begin{theorem}[A quadratic bound by an exponential]
\lean{Complexity.codec_square_le_pow}
\label{Complexity.codec_square_le_pow}
\leanok
\difficulty{5}
For every natural number $C \ge 4$, one has $C^2 \le 2^C$.
\end{theorem}

\begin{theorem}[Absorbing polynomial factors into an exponent]
\lean{Complexity.codec_absorb}
\label{Complexity.codec_absorb}
\leanok
\uses{Complexity.codec_square_le_pow}
\difficulty{6}
For any natural numbers $B,d,K$, there is a natural number $C>0$ with $K\le C$ such
that, for every natural number $u\ge2$, $B C^2 u^d\le C u^C$.
\end{theorem}

\begin{theorem}[Characterization of a position code]
\lean{Complexity.codec_posCode_eq_iff}
\label{Complexity.codec_posCode_eq_iff}
\leanok
\uses{Complexity.codec_posCode, Turing.MultiTapeTM.ConfigCount.posCode, Turing.MultiTapeTM.ConfigCount.posCode_injOn}
\difficulty{4}
Let $s$ be a natural number, let $z$ be an integer with $|z|\le s$, and let
$j\in\{0,1,\ldots,2s\}$. The position code of $z$ with bound $s$ equals $j$ if and only
if $z=j-s$.
\end{theorem}

\begin{definition}[Configuration bit observation]
\lean{Complexity.codecObserve}
\label{Complexity.codecObserve}
\leanok
\uses{Complexity.CodecBit, Complexity.codecRowBits, Complexity.codecRowOf, Turing.Cfg, Turing.FinTM}
Let $M$ be a finite Turing machine with Boolean alphabet, let $n,s$ be natural numbers,
let $x$ be a Boolean string of length $n$, and let $c$ be a configuration of $M$ on $x$.
For a work-tape coordinate $(i,j,b)$, the observed bit records whether cell $j-s$ on
tape $i$ contains false when $b=0$, contains true when $b=1$, or has the tape head at
$j-s$ when $b=2$. At every other codec-bit coordinate, it is the corresponding bit of
the row encoding of $c$.
\end{definition}

\begin{theorem}[Observation equals encoded row bit]
\lean{Complexity.codecObserve_eq}
\label{Complexity.codecObserve_eq}
\leanok
\uses{Complexity.codecObserve, Complexity.codecRowBits, Complexity.codecRowOf, Complexity.codec_posCode_eq_iff, Turing.Cfg, Turing.Cfg.InWindow, Turing.FinTM, Turing.MultiTapeTM.ConfigCount.coreCode}
\difficulty{4}
Let $M$ be a finite Turing machine over Booleans, let $n,s$ be natural numbers, let $x$
be a Boolean string of length $n$, and let $c$ be a configuration of $M$ on input $x$
whose work-tape heads lie in the window of radius $s$ and whose tapes are blank outside
that window. Then the Boolean observation of $c$ at every encoded bit position agrees
with the corresponding bit in the encoded row of $c$.
\end{theorem}

\begin{theorem}[Observation determined by core data]
\lean{Complexity.codecObserve_congr}
\label{Complexity.codecObserve_congr}
\leanok
\uses{Complexity.codecObserve, Complexity.codecRowBits, Complexity.codecRowOf, Turing.Cfg, Turing.FinTM, Turing.MultiTapeTM.ConfigCount.core, Turing.MultiTapeTM.ConfigCount.coreCode, Turing.NDTM.coreSum, Turing.outSummary}
\difficulty{4}
Let $M$ be a finite Turing machine over Booleans, let $n,s$ be natural numbers, let $x$
be a Boolean string of length $n$, and let $c,d$ be configurations of $M$ on input $x$.
If their core data together with their output summaries are equal, then their Boolean
observations agree at every encoded bit position.
\end{theorem}

\begin{theorem}[Blank cells remain blank after a step]
\lean{Complexity.codec_step_blank}
\label{Complexity.codec_step_blank}
\leanok
\uses{Turing.Action.apply, Turing.Cfg, Turing.Cfg.InWindow, Turing.Cfg.inputSymbol, Turing.Cfg.workTapeSymbols, Turing.FinTM, Turing.MultiTapeTM.step}
\difficulty{5}
Let $M$ be a finite Turing machine with Boolean work-tape symbols, let $x$ be a finite
Boolean input, and let $c$ be a configuration of $M$ on input $x$. Suppose that for some
natural number $s$, every work-tape head of $c$ is at an integer position of absolute
value at most $s$, and every work-tape cell at an integer position of absolute value
greater than $s$ is blank. Then, after one step of $M$ from $c$, every work-tape cell at
an integer position of absolute value greater than $s$ is blank.
\end{theorem}

\begin{theorem}[Step observation characterizes core equality]
\lean{Complexity.codecObserve_step_iff}
\label{Complexity.codecObserve_step_iff}
\leanok
\uses{Complexity.codecObserve, Complexity.codecObserve_congr, Complexity.codecObserve_eq, Complexity.codecRowBits, Complexity.codecRowBits_injective, Complexity.codecRowOf, Complexity.codecRowOf_eq_iff, Complexity.codec_step_blank, Turing.Cfg, Turing.Cfg.InWindow, Turing.FinTM, Turing.MultiTapeTM.step, Turing.NDTM.coreSum}
\difficulty{6}
Let $M$ be a finite Turing machine over Booleans, let $n,s$ be natural numbers, let $x$
be a Boolean string of length $n$, and let $c,d$ be configurations of $M$ on input $x$,
each with work-tape heads in the window of radius $s$ and tapes blank outside that
window. The observation of the one-step successor of $c$ agrees at every encoded bit
position with the encoded row of $d$ if and only if the successor of $c$ and $d$ have
equal core data and output summaries.
\end{theorem}

\begin{definition}[Output summary word]
\lean{Complexity.codecSummaryWord}
\label{Complexity.codecSummaryWord}
\leanok
\uses{Turing.OutSummary}
The output-summary word is the empty string for the empty-output summary, the one-bit
string $[\mathrm{true}]$ for the accepting summary, and the one-bit string
$[\mathrm{false}]$ for the dead summary.
\end{definition}

\begin{theorem}[Recovery of an output summary from its word]
\lean{Complexity.codecSummaryWord_summary}
\label{Complexity.codecSummaryWord_summary}
\leanok
\uses{Complexity.codecSummaryWord, Turing.OutSummary, Turing.outSummary}
\difficulty{1}
For every output summary $o$, the summary assigned to its encoded Boolean word is $o$.
\end{theorem}

\begin{theorem}[Step invariance under output normalization]
\lean{Complexity.codecObserve_step_normalize}
\label{Complexity.codecObserve_step_normalize}
\leanok
\uses{Complexity.codecObserve, Complexity.codecObserve_congr, Complexity.codecSummaryWord, Complexity.codecSummaryWord_summary, Turing.Cfg, Turing.FinTM, Turing.MultiTapeTM.step, Turing.MultiTapeTM.toNDTM, Turing.NDTM.coreSum_stepWith, Turing.outSummary}
\difficulty{3}
Let $M$ be a deterministic Turing machine with Boolean symbols, $n$ input positions, and
work-tape radius parameter $s$. For any Boolean input $x$ of length $n$ and
configuration $c$ on $x$, the observations of the one-step successors agree when the
output of $c$ is replaced by its summary word:
$\,\mathrm{observe}(M,n,s,x,\mathrm{step}(c))=\mathrm{observe}(M,n,s,x,\mathrm{step}(c'))$,
where $c'$ agrees with $c$ except that its output is the encoding of the summary of
$c$’s output.
\end{theorem}

\begin{definition}[Machine codec case]
\lean{Complexity.CodecCase}
\label{Complexity.CodecCase}
\leanok
\uses{Turing.FinTM}
For a finite Turing machine $M$ with Boolean alphabet, a codec case consists of an
optional state of $M$, an optional Boolean input symbol, and, for each work tape of $M$,
an optional Boolean symbol.
\end{definition}

\begin{definition}[Codec transition action]
\lean{Complexity.codecAction}
\label{Complexity.codecAction}
\leanok
\uses{Complexity.CodecCase, Turing.Action, Turing.FinTM}
Given a finite Turing machine $M$ with Boolean alphabet and a codec case consisting of
an optional state, an optional input symbol, and one optional symbol for each work tape,
define its action as follows: if the state is absent, the machine halts with no output,
no head movement, and no work-tape writes; if the state is present, the action is the
machine's transition on that state and the supplied symbols.
\end{definition}

\begin{theorem}[Encoded action reproduces a machine step]
\lean{Complexity.codecAction_apply}
\label{Complexity.codecAction_apply}
\leanok
\uses{Complexity.codecAction, Turing.Action.apply, Turing.Cfg, Turing.Cfg.inputSymbol, Turing.Cfg.workTapeSymbols, Turing.FinTM, Turing.MultiTapeTM.step}
\difficulty{4}
For every finite Turing machine $M$, input word $x$, and configuration $c$ on input $x$,
applying the action encoded from $c$—using the machine's transition action when the
state is present and a stationary, nonwriting halting action otherwise—gives exactly the
configuration obtained by one step of $M$.
\end{theorem}

\begin{definition}[Codec port labels]
\lean{Complexity.CodecPort}
\label{Complexity.CodecPort}
\leanok
\uses{Turing.FinTM}
For a finite Turing machine $M$ with Boolean alphabet, the codec ports consist of labels
for the state, input head, and input symbol; for each work tape, a head label and two
symbol labels, one for each Boolean value; two summary labels, one for each Boolean
value; and labels for the old and new bits.
\end{definition}

\begin{theorem}[Number of codec ports]
\lean{Complexity.codecPort_card}
\label{Complexity.codecPort_card}
\leanok
\uses{Complexity.CodecPort, Turing.FinTM}
\difficulty{3}
For every finite-state Turing machine $M$ with Boolean symbols, the number of codec
ports is $3k+7$, where $k$ is the number of work tapes of $M$.
\end{theorem}

\begin{definition}[Boolean symbol decoding]
\lean{Complexity.codecReadSymbol}
\label{Complexity.codecReadSymbol}
\leanok
For Boolean inputs $f$ and $t$, define the decoded optional symbol to be true if $t$ is
true, false if $t$ is false and $f$ is true, and absent if both are false.
\end{definition}

\begin{theorem}[Recovery of a Boolean symbol]
\lean{Complexity.codecReadSymbol_flags}
\label{Complexity.codecReadSymbol_flags}
\leanok
\uses{Complexity.codecReadSymbol}
\difficulty{4}
For every optional Boolean symbol $v$, reading the flags that indicate whether $v$ is
respectively the present symbol false or the present symbol true recovers $v$.
\end{theorem}

\begin{definition}[Output summary decoding]
\lean{Complexity.codecReadSummary}
\label{Complexity.codecReadSummary}
\leanok
\uses{Turing.OutSummary}
For Boolean inputs $acc$ and $dead$, define the output summary to be dead if $dead$ is
true, accept if $dead$ is false and $acc$ is true, and empty if both are false.
\end{definition}

\begin{theorem}[Recovery of an output summary]
\lean{Complexity.codecReadSummary_flags}
\label{Complexity.codecReadSummary_flags}
\leanok
\uses{Complexity.codecReadSummary, Turing.OutSummary}
\difficulty{4}
For every output summary $o$, reading the two flags that indicate whether $o$ is the
accepting summary or the dead summary recovers $o$.
\end{theorem}

\begin{definition}[Input wire encoding]
\lean{Complexity.codecInputWire}
\label{Complexity.codecInputWire}
\leanok
\uses{Complexity.CodecBit, Turing.FinTM}
Let $M$ be a finite Turing machine with Boolean alphabet, and let $n,s$ be natural
numbers. For a position $p\in\{0,\ldots,n+1\}$, encode it by the input bit with index
$p-1$ when $1\le p\le n$; otherwise encode it by the summary bit indexed by false.
\end{definition}

\begin{definition}[Input symbol reading]
\lean{Complexity.codecInputRead}
\label{Complexity.codecInputRead}
\leanok
Let $n$ be a natural number, let $p$ be one of the $n+2$ positions, and let $b$ be a
Boolean. The decoded input symbol at position $p$ is the optional symbol $b$ when $1\le
p\le n$, and is absent otherwise.
\end{definition}

\begin{definition}[Codec gate wire mapping]
\lean{Complexity.codecGateWire}
\label{Complexity.codecGateWire}
\leanok
\uses{Complexity.CodecBit, Complexity.CodecCase, Complexity.CodecPort, Complexity.CodecPositions, Complexity.codecInputWire, Turing.FinTM}
Let $M$ be a finite Turing machine with Boolean alphabet, let $n,s$ be natural numbers,
let $p$ specify an input-head position and a work-head position for each work tape, let
$q$ be a codec case, and let $b$ be a codec bit. The gate-wire map assigns each port a
Boolean flag and a codec bit: the state, input-head, input-symbol, work-head,
work-symbol, and summary ports all have flag false and carry, respectively, the
corresponding state bit, head bit, input-wire bit, work-head bit, work-symbol bit (using
index $1$ for true and $0$ for false), or summary bit. The old-bit port has flag false
and carries $b$, while the new-bit port has flag true and carries $b$.
\end{definition}

\begin{definition}[Codec configuration guard]
\lean{Complexity.codecGuard}
\label{Complexity.codecGuard}
\leanok
\uses{Complexity.CodecCase, Complexity.CodecPort, Complexity.CodecPositions, Complexity.codecInputRead, Complexity.codecReadSymbol, Turing.FinTM}
Let $M$ be a finite Turing machine with Boolean alphabet, let $n,s$ be natural numbers,
let $p$ specify an input-head position and a work-head position for each work tape, let
$q$ be a codec case, and let $a$ assign a Boolean to each codec port. The assignment
satisfies the guard for $p$ and $q$ precisely when the state, input-head, and every
work-head port are true; reading the input-symbol port at the input-head position gives
the input symbol in $q$; and, for each work tape, reading its two symbol ports gives the
work symbol in $q$.
\end{definition}

\begin{theorem}[Input symbol read from its wire]
\lean{Complexity.codecInputRead_wire}
\label{Complexity.codecInputRead_wire}
\leanok
\uses{Complexity.codecInputRead, Complexity.codecInputWire, Complexity.codecRowBits, Complexity.codecRowOf, Turing.Cfg, Turing.Cfg.inputSymbol, Turing.FinTM}
\difficulty{5}
Let $M$ be a finite Turing machine over Booleans, let $n,s$ be natural numbers, let $x$
be a Boolean string of length $n$, and let $c$ be a configuration of $M$ on input $x$.
For any input-head position $p$, if the input-head position in the encoding of $c$ is
$p$, then the Boolean symbol obtained by reading the encoded input wire at $p$ (with
boundary positions yielding no symbol) equals the symbol scanned by $c$ on its input
tape.
\end{theorem}

\begin{theorem}[Guard condition in row form]
\lean{Complexity.codecGuard_rows}
\label{Complexity.codecGuard_rows}
\leanok
\uses{Complexity.CodecBit, Complexity.CodecCase, Complexity.CodecPositions, Complexity.CodecRow, Complexity.codecGateWire, Complexity.codecGuard, Complexity.codecInputRead, Complexity.codecInputWire, Complexity.codecPairBits, Complexity.codecReadSymbol_flags, Complexity.codecRowBits, Turing.FinTM}
\difficulty{1}
Let $M$ be a finite Turing machine with $k$ work tapes, let $n,s$ be natural numbers,
let $p$ be a tuple of an input-head position and work-head positions, let $q$ be a case
consisting of an optional state, an optional input symbol, and one optional symbol per
work tape, let $b$ be a codec bit, and let $r,t$ be codec rows. The wired port
assignment satisfies the codec guard if and only if the state and input-head position in
$r$ equal those in $q$ and $p$, respectively; each work-tape head in $r$ equals its
position in $p$; the encoded input symbol read at $p$ equals the input symbol in $q$;
and, for every work tape, the symbol recorded in $r$ at its position in $p$ equals the
corresponding symbol in $q$.
\end{theorem}

\begin{theorem}[Configuration decoded by a guard]
\lean{Complexity.codecGuard_spec}
\label{Complexity.codecGuard_spec}
\leanok
\uses{Complexity.CodecBit, Complexity.CodecCase, Complexity.CodecPositions, Complexity.CodecRow, Complexity.codecGateWire, Complexity.codecGuard, Complexity.codecGuard_rows, Complexity.codecInputRead_wire, Complexity.codecPairBits, Complexity.codecRowOf, Complexity.codec_posCode_eq_iff, Turing.Cfg, Turing.Cfg.InWindow, Turing.Cfg.inputSymbol, Turing.Cfg.workTapeSymbols, Turing.FinTM, Turing.MultiTapeTM.ConfigCount.coreCode}
\difficulty{4}
Let $M$ be a finite Turing machine over the Boolean alphabet, let $n,s$ be natural
numbers, and let $x$ be a Boolean string of length $n$. Let $c$ be a configuration of
$M$ on input $x$ whose work-tape heads lie in the window of radius $s$ and whose work
tapes are blank outside that window. For any encoded row $t$, tuple of input and
work-tape positions $p$, case $q$ specifying a state and scanned symbols, and bit index
$b$, the guard formed by reading the wires selected by $p,q,b$ from the pair consisting
of the encoding of $c$ and $t$ holds exactly when $c$ has state $q$'s state, its encoded
input-head position is $p$'s input position, each work-tape head is at position
$(p_i-s)$, and its scanned input and work-tape symbols are those specified by $q$.
\end{theorem}

\begin{definition}[Expected Configuration Bits]
\lean{Complexity.codecExpected}
\label{Complexity.codecExpected}
\leanok
\uses{Complexity.CodecBit, Complexity.CodecPort, Complexity.CodecPositions, Complexity.codecReadSummary, Complexity.codecSummaryWord, Turing.Action, Turing.FinTM, Turing.moveInputPos, Turing.outSummary}
Let $M$ be a finite Turing machine with Boolean input alphabet, let $n,s\in\mathbb{N}$,
and choose a bounded input-head position and a bounded work-head position for each tape.
Given a machine action and a Boolean assignment to the codec ports, the expected value
of each encoded bit is defined as follows: input bits copy the old bit; the input-head
and state bits indicate whether their indices or states match the action's resulting
position or successor state; work-symbol bits copy the old bit except at the scanned
cell, where they indicate whether the action writes that symbol; work-head bits indicate
whether the head moves to that cell; and summary bits indicate whether the updated
output summary is accepting or dead.
\end{definition}

\begin{definition}[Codec Gate Predicate]
\lean{Complexity.codecGate}
\label{Complexity.codecGate}
\leanok
\uses{Complexity.CodecBit, Complexity.CodecCase, Complexity.CodecPort, Complexity.CodecPositions, Complexity.codecAction, Complexity.codecExpected, Complexity.codecGuard, Turing.FinTM}
Let $M$ be a finite Turing machine with Boolean input alphabet, let $n,s\in\mathbb{N}$,
and choose codec positions, a codec case, an encoded bit, and a Boolean assignment to
the codec ports. The gate is true exactly when either the assignment does not satisfy
the codec guard for these data, or its new-bit value equals the expected value of the
chosen encoded bit under the action determined by the codec case.
\end{definition}

\begin{theorem}[Expected encoding of an action]
\lean{Complexity.codecExpected_apply}
\label{Complexity.codecExpected_apply}
\leanok
\uses{Complexity.CodecBit, Complexity.CodecCase, Complexity.CodecPositions, Complexity.CodecRow, Complexity.codecExpected, Complexity.codecGateWire, Complexity.codecObserve, Complexity.codecPairBits, Complexity.codecReadSummary_flags, Complexity.codecRowBits, Complexity.codecRowOf, Complexity.codecSummaryWord, Turing.Action, Turing.Action.apply, Turing.Cfg, Turing.FinTM, Turing.MultiTapeTM.ConfigCount.coreCode, Turing.moveInputPos, Turing.outSummary}
\difficulty{5}
Let $M$ be a finite-state Turing machine with $k$ work tapes, let $n,s$ be natural
numbers, and let $x$ be a Boolean input of length $n$. Let $c$ be a configuration on
input $x$, $t$ a codec row, $p$ a tuple of an input-head position and work-head
positions, $q$ a case consisting of an optional state, an optional input symbol, and one
optional symbol per work tape, and $a$ a Boolean assignment to the codec ports. Assume
that the input head recorded in the row of $c$ is the input-head position in $p$, and
that each work-tape head of $c$ is at the integer position represented by its
corresponding position in $p$ (with offset $s$). For every codec bit $b$, the expected
value of $b$ under the action determined by $q$, using the port values read from the row
of $c$ and $t$ by the gate wiring, equals the observed value of $b$ in the configuration
obtained by applying that action to $c$ after replacing its output by the encoding of
its output summary.
\end{theorem}

\begin{theorem}[Expected encoding of a machine step]
\lean{Complexity.codecExpected_step}
\label{Complexity.codecExpected_step}
\leanok
\uses{Complexity.CodecBit, Complexity.CodecCase, Complexity.CodecPositions, Complexity.CodecRow, Complexity.codecAction, Complexity.codecAction_apply, Complexity.codecExpected, Complexity.codecExpected_apply, Complexity.codecGateWire, Complexity.codecGuard, Complexity.codecGuard_spec, Complexity.codecObserve, Complexity.codecObserve_step_normalize, Complexity.codecPairBits, Complexity.codecRowOf, Complexity.codecSummaryWord, Turing.Action.apply, Turing.Cfg, Turing.Cfg.InWindow, Turing.Cfg.inputSymbol, Turing.Cfg.workTapeSymbols, Turing.FinTM, Turing.MultiTapeTM.step, Turing.outSummary}
\difficulty{4}
Let $M$ be a finite-state Turing machine with $k$ work tapes, let $n,s$ be natural
numbers, and let $x$ be a Boolean input of length $n$. Let $c$ be a configuration on
input $x$ whose work-tape heads lie within distance $s$ of the origin and whose work
tapes are blank outside that window. Let $t$ be a codec row, $p$ a tuple of an
input-head position and work-head positions, $q$ a case consisting of an optional state,
an optional input symbol, and one optional symbol per work tape, and $a$ a Boolean
assignment to the codec ports. If the wired port values satisfy the codec guard for $p$
and $q$, then for every codec bit $b$, its expected value under the action determined by
$q$ equals its observed value in the successor configuration of $c$ under one step of
$M$.
\end{theorem}

\begin{theorem}[Existence of a gate CNF]
\lean{Complexity.codecGate_cnf}
\label{Complexity.codecGate_cnf}
\leanok
\uses{Complexity.CodecBit, Complexity.CodecCase, Complexity.CodecPort, Complexity.CodecPositions, Complexity.codecCnfCost, Complexity.codecGate, Complexity.codecGateWire, Complexity.codecLocalCheck, Complexity.codecPort_card, Turing.FinTM}
\difficulty{3}
Let $M$ be a finite Turing machine with $k$ work tapes, let $n,s$ be natural numbers,
let $p$ be a tuple of an input-head position and work-head positions, let $q$ be a case
consisting of an optional state, an optional input symbol, and one optional symbol per
work tape, and let $b$ be a codec bit. There exists a CNF formula whose cost is at most
$2^{3k+7}(3k+8)$ and which, under every Boolean assignment to its variables, evaluates
to the gate value obtained by wiring those variables to the prescribed ports.
\end{theorem}

\begin{definition}[CNF for a Codec Gate]
\lean{Complexity.codecGateCNF}
\label{Complexity.codecGateCNF}
\leanok
\uses{Complexity.CodecBit, Complexity.CodecCase, Complexity.CodecPositions, Complexity.codecGate_cnf, Turing.FinTM}
Let $M$ be a finite Turing machine with Boolean input alphabet, let $n,s\in\mathbb{N}$,
and choose codec positions, a codec case, and an encoded bit. The associated gate
formula is represented by the CNF selected from the gate's CNF encoding, with variables
indexed by a Boolean value and an encoded bit.
\end{definition}

\begin{theorem}[Specification of the gate CNF]
\lean{Complexity.codecGateCNF_spec}
\label{Complexity.codecGateCNF_spec}
\leanok
\uses{Complexity.CodecBit, Complexity.CodecCase, Complexity.CodecPositions, Complexity.codecCnfCost, Complexity.codecGate, Complexity.codecGateCNF, Complexity.codecGateWire, Complexity.codecGate_cnf, Turing.FinTM}
\difficulty{1}
Let $M$ be a finite Turing machine with $k$ work tapes, let $n,s$ be natural numbers,
let $p$ be a tuple of an input-head position and work-head positions, let $q$ be a case
consisting of an optional state, an optional input symbol, and one optional symbol per
work tape, and let $b$ be a codec bit. The cost of the CNF formula assigned to these
data is at most $2^{3k+7}(3k+8)$, and for every Boolean assignment $a$ to its variables,
the formula evaluates to the gate value on the port assignment obtained by reading the
wired variables from $a$.
\end{theorem}

\begin{definition}[Codec Gate Indices]
\lean{Complexity.CodecGateIndex}
\label{Complexity.CodecGateIndex}
\leanok
\uses{Complexity.CodecBit, Complexity.CodecCase, Complexity.CodecPositions, Turing.FinTM}
For a finite Turing machine $M$ with Boolean input alphabet and natural numbers $n,s$, a
codec gate index is a triple consisting of codec positions, a codec case, and an encoded
bit.
\end{definition}

\begin{definition}[CNF for Machine Transitions]
\lean{Complexity.codecTransitionCNF}
\label{Complexity.codecTransitionCNF}
\leanok
\uses{Complexity.CodecBit, Complexity.CodecGateIndex, Complexity.codecGateCNF, Turing.FinTM}
Let $M$ be a finite Turing machine with Boolean input alphabet, and let
$n,s\in\mathbb{N}$. The transition CNF is the concatenation, over all codec gate
indices, of the CNFs associated with the corresponding gates.
\end{definition}

\begin{theorem}[Transition Formula Evaluation]
\lean{Complexity.codecTransitionCNF_eval}
\label{Complexity.codecTransitionCNF_eval}
\leanok
\uses{Complexity.CodecBit, Complexity.CodecGateIndex, Complexity.codecCnf_eval_flatMap, Complexity.codecGate, Complexity.codecGateCNF_spec, Complexity.codecGateWire, Complexity.codecTransitionCNF, Turing.FinTM}
\difficulty{3}
Let $M$ be a finite-state Turing machine with $k$ work tapes, let $n,s\in\mathbb{N}$,
and assign a Boolean value to every encoded bit. The transition CNF evaluates to true
exactly when, for every choice of a position, a transition case, and an encoded bit, the
corresponding gate evaluates to true on the assignment induced by its gate wires.
\end{theorem}

\begin{theorem}[Transition Formula Characterizes a Step]
\lean{Complexity.codecTransitionCNF_step_iff}
\label{Complexity.codecTransitionCNF_step_iff}
\leanok
\uses{Complexity.CodecCase, Complexity.CodecPositions, Complexity.codecExpected_step, Complexity.codecGate, Complexity.codecGateWire, Complexity.codecGuard, Complexity.codecGuard_spec, Complexity.codecObserve_step_iff, Complexity.codecPairBits, Complexity.codecRowBits, Complexity.codecRowOf, Complexity.codecTransitionCNF, Complexity.codecTransitionCNF_eval, Complexity.codec_posCode_eq_iff, Turing.Cfg, Turing.Cfg.InWindow, Turing.Cfg.inputSymbol, Turing.Cfg.workTapeSymbols, Turing.FinTM, Turing.MultiTapeTM.step, Turing.NDTM.coreSum}
\difficulty{5}
Let $M$ be a finite-state Turing machine with $k$ work tapes, let $n,s\in\mathbb{N}$,
and let $x$ be a Boolean input of length $n$. For configurations $c,d$ on input $x$,
suppose each has all work-tape head positions within distance $s$ of the origin and has
blank contents outside that window. The transition CNF evaluated on the pair of their
encoded bit assignments is true exactly when the core and output summary of the one-step
successor of $c$ equal those of $d$.
\end{theorem}

\begin{theorem}[Transition Formula Cost]
\lean{Complexity.codecTransitionCNF_cost}
\label{Complexity.codecTransitionCNF_cost}
\leanok
\uses{Complexity.CodecBit, Complexity.CodecCase, Complexity.CodecGateIndex, Complexity.CodecPositions, Complexity.codecCnfCost, Complexity.codecCnfCost_flatMap_le, Complexity.codecGateCNF, Complexity.codecGateCNF_spec, Complexity.codecTransitionCNF, Turing.FinTM}
\difficulty{3}
Let $M$ be a finite-state Turing machine with $k$ work tapes, and let
$n,s\in\mathbb{N}$. The cost of the transition CNF is at most the product of the numbers
of position choices, transition cases, and encoded bits, multiplied by $2^{3k+7}(3k+8)$.
\end{theorem}

\begin{definition}[Configuration Adjacency Formula]
\lean{Complexity.configurationAdjacency}
\label{Complexity.configurationAdjacency}
\leanok
\uses{Complexity.codecInputEquality, Complexity.codecPairWire, Complexity.codecTransitionCNF, Turing.FinTM}
Let $M$ be a finite Turing machine with Boolean input alphabet, and let
$C,n,s\in\mathbb{N}$. The configuration-adjacency CNF is obtained by concatenating the
input-equality CNF with the transition CNF, then relabelling its variables by the
specified pairing of wires.
\end{definition}

\begin{theorem}[Evaluation of the Adjacency Formula]
\lean{Complexity.configurationAdjacency_eval}
\label{Complexity.configurationAdjacency_eval}
\leanok
\uses{Complexity.codecCoefficient, Complexity.codecInputEquality_eval, Complexity.codecTransitionCNF_step_iff, Complexity.configurationAdjacency, Complexity.configurationCodec, Complexity.configurationCodec_pair_get, Turing.Cfg, Turing.Cfg.InWindow, Turing.FinTM, Turing.MultiTapeTM.step, Turing.NDTM.coreSum}
\difficulty{3}
Let $M$ be a finite-state Turing machine with $k$ work tapes, and let $n,s,C\in\mathbb
N$ with $C\ge |Q|+6k+5$, where $Q$ is its state set. Let $x$ be a Boolean input of
length $n$, and let $c,d$ be configurations of $M$ on $x$ whose work-tape heads lie
within distance $s$ of the origin and whose work tapes are blank outside that interval.
Evaluating the adjacency formula on the concatenation of the encodings of $c$ and $d$
gives true exactly when the core data of the successor of $c$ under one machine step,
together with its output summary, equals the corresponding data of $d$.
\end{theorem}

\begin{theorem}[Rejection for Different Inputs]
\lean{Complexity.configurationAdjacency_reject}
\label{Complexity.configurationAdjacency_reject}
\leanok
\uses{Complexity.codecCoefficient, Complexity.configurationAdjacency, Complexity.configurationCodec, Complexity.configurationInputEquality, Complexity.configurationInputEquality_reject, Turing.Cfg, Turing.FinTM}
\difficulty{3}
Let $M$ be a finite-state Turing machine with $k$ work tapes, and let $n,s,C\in\mathbb
N$ with $C\ge |Q|+6k+5$, where $Q$ is its state set. If $x$ and $y$ are distinct Boolean
inputs, each of length $n$, then for any configurations $c$ on $x$ and $d$ on $y$, the
adjacency formula evaluated on the concatenation of their configuration encodings is
false.
\end{theorem}

\begin{definition}[Transition Encoding Coefficient]
\lean{Complexity.codecTransitionCoefficient}
\label{Complexity.codecTransitionCoefficient}
\leanok
\uses{Complexity.CodecCase, Turing.FinTM}
For a finite Turing machine $M$ with Boolean input alphabet and $k$ work tapes, define
the transition encoding coefficient by $2^k\,|\mathcal{Q}|\,2^{3k+7}(3k+8)$, where
$\mathcal{Q}$ is the set of codec cases for $M$.
\end{definition}

\begin{theorem}[Polynomial Transition Formula Cost]
\lean{Complexity.codecTransitionCNF_cost_poly}
\label{Complexity.codecTransitionCNF_cost_poly}
\leanok
\uses{Complexity.CodecBit, Complexity.CodecCase, Complexity.CodecPositions, Complexity.codecBit_card_le, Complexity.codecCnfCost, Complexity.codecCoefficient, Complexity.codecPositions_card_le, Complexity.codecTransitionCNF, Complexity.codecTransitionCNF_cost, Complexity.codecTransitionCoefficient, Turing.FinTM}
\difficulty{4}
Let $M$ be a finite-state Turing machine with $k$ work tapes, and let
$C,n,s\in\mathbb{N}$. If the coefficient $|\mathrm{State}(M)|+6k+5$ is at most $C$, then
the transition CNF cost is at most $2^k$ times the number of transition cases, times
$2^{3k+7}(3k+8)$, times $C(s+n+2)^{k+2}$.
\end{theorem}

\begin{theorem}[Adjacency encoding cost bound]
\lean{Complexity.codecAdjacency_cost}
\label{Complexity.codecAdjacency_cost}
\leanok
\uses{Complexity.codecCnfCost, Complexity.codecCnfCost_append, Complexity.codecCoefficient, Complexity.codecInputEquality, Complexity.codecInputEquality_cost, Complexity.codecTransitionCNF, Complexity.codecTransitionCNF_cost_poly, Complexity.codecTransitionCoefficient, Turing.FinTM}
\difficulty{4}
Let $M$ be a finite-state Turing machine over the Boolean alphabet, with $k$ work tapes,
and let $n,s,C$ be nonnegative integers. Write $Q$ for its state set, and let
$A=|Q|+6k+5$ and $T=2^k(2|Q|+2)^{	ext{?}}$—more precisely,
$T=2^k\,|\mathrm{Option}(Q)\times\mathrm{Option}(\mathrm{Bool})\times(\mathrm{Fin}(k)\to\mathrm{Option}(\mathrm{Bool}))|\,2^{3k+7}(3k+8)$.
If $A\le C$, then the cost of the concatenation of the input-equality CNF and transition
CNF is at most $(6+T)C(s+n+2)^{k+2}$. Here the cost of a CNF is the sum, over its
clauses, of one plus the clause length.
\end{theorem}

\begin{theorem}[Configuration Adjacency Bounds]
\lean{Complexity.configurationAdjacency_bounds}
\label{Complexity.configurationAdjacency_bounds}
\leanok
\uses{Complexity.codecAdjacency_cost, Complexity.codecCnfCost, Complexity.codecCnf_numVars, Complexity.codecCnf_serialize, Complexity.codecCoefficient, Complexity.codecInputEquality, Complexity.codecPairWire, Complexity.codecPairWire_lt, Complexity.codecTransitionCNF, Complexity.codecTransitionCoefficient, Complexity.configurationAdjacency, Std.Sat.CNF.numVars, Std.Sat.CNF.serialize, Turing.FinTM}
\difficulty{4}
Let $M$ be a finite-state Turing machine with $k$ work tapes, and let $n,s,C\in\mathbb
N$ with $C\ge |Q|+6k+5$, where $Q$ is its state set. The adjacency formula has variable
count at most $2C(s+n+1)$, and its binary serialization has length at most
$\bigl(5(6+D)+1\bigr)C^2(s+n+2)^{k+3}$, where $D=2^k\,|\mathcal Q|\,2^{3k+7}(3k+8)$ and
$\mathcal Q$ is the set of possible transition cases.
\end{theorem}

\begin{theorem}[CNF Encodings of Machine Configurations]
\lean{Complexity.exists_adjacency_codec_cnf}
\label{Complexity.exists_adjacency_codec_cnf}
\leanok
\uses{Complexity.codecCoefficient, Complexity.codecTransitionCoefficient, Complexity.codecValidityCoefficient, Complexity.codec_absorb, Complexity.configurationAcceptance, Complexity.configurationAcceptance_bounds, Complexity.configurationAcceptance_eval, Complexity.configurationAcceptance_size, Complexity.configurationAdjacency, Complexity.configurationAdjacency_bounds, Complexity.configurationAdjacency_eval, Complexity.configurationAdjacency_reject, Complexity.configurationCodec, Complexity.configurationCodec_eq_iff, Complexity.configurationCodec_input_eq, Complexity.configurationCodec_length, Complexity.configurationValidity, Complexity.configurationValidity_eval, Complexity.configurationValidity_numVars, Complexity.configurationValidity_size, Std.Sat.CNF.numVars, Std.Sat.CNF.serialize, Turing.Cfg, Turing.Cfg.InWindow, Turing.FinTM, Turing.MultiTapeTM.step, Turing.NDTM.coreSum, Turing.outSummary}
\difficulty{5}
Let $M$ be a deterministic Turing machine over the binary alphabet, with a finite state
set, and let $n,s$ be natural numbers. There is a positive integer $C$, independent of
$n,s$, such that for every $n,s$ one can choose an encoding of configurations and CNF
formulas $\varphi_v,\varphi_a,\varphi_{\mathrm{acc}}$ with these properties: every
configuration encoding has length $C(s+n+1)$; for inputs of length $n$ whose work-tape
heads and nonblank cells lie within distance $s$ of the origin, equal encodings imply
equal inputs, and for a fixed input two configurations have equal encodings exactly when
their states, input-head positions, work tapes, work-tape-head positions, and output
summaries agree. The formula $\varphi_v$, evaluated on any bit string of that encoding
length (with variables beyond the string assigned false), accepts exactly the encodings
of such configurations on inputs of length $n$. The formula $\varphi_a$, evaluated on
the concatenation of two configuration encodings, accepts exactly when the first
configuration's one-step successor under $M$ has the same state, input-head position,
work tapes, and work-tape-head positions as the second; it rejects every such pair whose
inputs differ. The formula $\varphi_{\mathrm{acc}}$ accepts a configuration encoding
exactly when the configuration is halted and its output is precisely the one-bit string
$1$. The variable counts of $\varphi_v$ and $\varphi_{\mathrm{acc}}$ are at most
$C(s+n+1)$, and that of $\varphi_a$ is at most $2C(s+n+1)$. Each of their binary CNF
encodings has length at most $C(s+n+2)^C$.
\end{theorem}

\begin{theorem}[Exact configuration width]
\lean{Complexity.ConfigCodec.width_eq}
\label{Complexity.ConfigCodec.width_eq}
\leanok
\uses{Complexity.ConfigCodec.stateCount, Complexity.ConfigCodec.stateOffset, Complexity.ConfigCodec.summaryOffset, Complexity.ConfigCodec.width, Complexity.ConfigCodec.workOffset}
\difficulty{3}
Let $M$ be a finite Turing machine over the Boolean alphabet, with $k$ work tapes, and
let $n,s\in\mathbb{N}$. If $W_M(n,s)$ denotes the width obtained by the configuration
encoding, then $W_M(n,s)=2n+4+|\mathrm{State}(M)|+1+3k(2s+1)$, where
$|\mathrm{State}(M)|+1$ is the number of machine states plus one.
\end{theorem}

\begin{theorem}[Configuration width bound]
\lean{Complexity.ConfigCodec.width_le}
\label{Complexity.ConfigCodec.width_le}
\leanok
\uses{Complexity.ConfigCodec.width}
Let $M$ be a finite Turing machine over the Boolean alphabet, with $k$ work tapes, and
let $n,s\in\mathbb{N}$. The width $W_M(n,s)$ of its configuration encoding satisfies
$W_M(n,s)\le (|\mathrm{State}(M)|+6k+5)(s+n+1)$, where $|\mathrm{State}(M)|$ is the
number of machine states.
\end{theorem}

\begin{theorem}[Index of the Halting State Variable]
\lean{Complexity.ConfigCodec.stateVar_none}
\label{Complexity.ConfigCodec.stateVar_none}
\leanok
\uses{Complexity.ConfigCodec.stateIndex, Complexity.ConfigCodec.stateOffset, Complexity.ConfigCodec.stateVar}
\difficulty{1}
For every finite-state Turing machine $M$ with Boolean tape symbols and every natural
number $n$, the variable assigned to the absent state (the halting state) is the state
offset $2n+2$.
\end{theorem}

\begin{theorem}[Cardinality of the port set]
\lean{Complexity.ConfigCodec.card_port}
\label{Complexity.ConfigCodec.card_port}
\leanok
\uses{Complexity.ConfigCodec.Port, Complexity.ConfigCodec.portCount}
\difficulty{3}
For every $k\in\mathbb{N}$, the set of ports for $k$ work tapes has cardinality $3k+7$.
\end{theorem}

\begin{theorem}[Port Indices Fit Within the Port Count]
\lean{Complexity.ConfigCodec.portIndex_lt}
\label{Complexity.ConfigCodec.portIndex_lt}
\leanok
\uses{Complexity.ConfigCodec.Port, Complexity.ConfigCodec.portCount, Complexity.ConfigCodec.portIndex}
\difficulty{5}
For every natural number $k$ and every port $p$ among the state, input-head,
input-symbol, work-head, work-symbol, summary, old-bit, and new-bit ports for $k$ work
tapes, the index assigned to $p$ is strictly less than $3k+7$.
\end{theorem}

\begin{theorem}[Port Decoding Recovers the Port]
\lean{Complexity.ConfigCodec.portOf_portIndex}
\label{Complexity.ConfigCodec.portOf_portIndex}
\leanok
\uses{Complexity.ConfigCodec.Port, Complexity.ConfigCodec.portIndex, Complexity.ConfigCodec.portOf}
\difficulty{6}
For every natural number $k$ and every port $p$ among the state, input-head,
input-symbol, work-head, work-symbol, summary, old-bit, and new-bit ports for $k$ work
tapes, decoding the index assigned to $p$ returns $p$.
\end{theorem}

\begin{theorem}[Length of the configuration code]
\lean{Complexity.ConfigCodec.code_length}
\label{Complexity.ConfigCodec.code_length}
\leanok
\uses{Complexity.ConfigCodec.code, Complexity.ConfigCodec.width, Turing.Cfg}
Let $M$ be a finite Turing machine over the Boolean alphabet, let $n,s$ be natural
numbers, let $x$ be a Boolean string, and let $c$ be a configuration of $M$ on input
$x$. The length of the configuration code is its prescribed width.
\end{theorem}

\begin{theorem}[Input segment of the configuration code]
\lean{Complexity.ConfigCodec.code_getD_inputVar}
\label{Complexity.ConfigCodec.code_getD_inputVar}
\leanok
\uses{Complexity.ConfigCodec.code, Complexity.ConfigCodec.inputVar, Turing.Cfg}
Let $M$ be a finite Turing machine over the Boolean alphabet, let $n,s$ be natural
numbers, let $x$ be a Boolean string, and let $c$ be a configuration of $M$ on input
$x$. For every natural number $i<n$, the bit of the configuration code at position $i$
(taking the default value false if that position is out of range) equals the $i$-th bit
of $x$, again defaulting to false if $i$ is beyond the end of $x$.
\end{theorem}

\begin{theorem}[Input-head bit in the configuration code]
\lean{Complexity.ConfigCodec.code_getD_headVar}
\label{Complexity.ConfigCodec.code_getD_headVar}
\leanok
\uses{Complexity.ConfigCodec.code, Complexity.ConfigCodec.headVar, Turing.Cfg}
Let $M$ be a finite Turing machine with Boolean tape symbols, let $n,s\in\mathbb{N}$,
let $x$ be a Boolean input, and let $c$ be a configuration of $M$ on input $x$. For
every $j<n+2$, the bit of the code at position $n+j$ is true exactly when the input-head
position of $c$ is $j$.
\end{theorem}

\begin{theorem}[State indicators in the configuration code]
\lean{Complexity.ConfigCodec.code_getD_stateVar}
\label{Complexity.ConfigCodec.code_getD_stateVar}
\leanok
\uses{Complexity.ConfigCodec.code, Complexity.ConfigCodec.stateVar, Turing.Cfg}
Let $M$ be a finite Turing machine over the Boolean alphabet, let $n,s$ be natural
numbers, let $x$ be a Boolean string, let $c$ be a configuration of $M$ on input $x$,
and let $q$ be either a state of $M$ or the halting marker. The bit of the configuration
code at the position assigned to $q$ by the state encoding (with false as the default
out-of-range value) is true exactly when the current state of $c$ equals $q$.
\end{theorem}

\begin{theorem}[Work tape flags in the configuration code]
\lean{Complexity.ConfigCodec.code_getD_workVar}
\label{Complexity.ConfigCodec.code_getD_workVar}
\leanok
\uses{Complexity.ConfigCodec.code, Complexity.ConfigCodec.workFlags, Complexity.ConfigCodec.workVar, Turing.Cfg}
Let $M$ be a finite Turing machine over the Boolean alphabet, let $n,s$ be natural
numbers, let $x$ be a Boolean string, let $c$ be a configuration of $M$ on input $x$,
let $i$ index one of its work tapes, and let $j,b$ be natural numbers with $j<2s+1$ and
$b<3$. The bit of the configuration code at the position assigned to work tape $i$, cell
index $j$, and flag index $b$ (defaulting to false if out of range) equals the $b$-th
bit of the three work flags for that tape and cell, with false as the default if the
flag index is out of range.
\end{theorem}

\begin{theorem}[Output summary indicators in the configuration code]
\lean{Complexity.ConfigCodec.code_getD_summaryVar}
\label{Complexity.ConfigCodec.code_getD_summaryVar}
\leanok
\uses{Complexity.ConfigCodec.code, Complexity.ConfigCodec.summaryVar, Turing.Cfg, Turing.outSummary}
Let $M$ be a finite Turing machine over the Boolean alphabet, let $n,s$ be natural
numbers, let $x$ be a Boolean string, and let $c$ be a configuration of $M$ on input
$x$. For either choice of the dead indicator, the bit at the corresponding summary
position in the configuration code, defaulting to false if out of range, is true exactly
when the output summary of $c$ is dead for the dead indicator and accepting otherwise.
\end{theorem}

\begin{theorem}[Input determined by the configuration code]
\lean{Complexity.ConfigCodec.code_input_eq}
\label{Complexity.ConfigCodec.code_input_eq}
\leanok
\uses{Complexity.ConfigCodec.code, Turing.Cfg}
Let $M$ be a finite Turing machine over the Boolean alphabet, let $n,s$ be natural
numbers, and let $x,x'$ be Boolean strings each of length $n$. For configurations $c$ of
$M$ on input $x$ and $c'$ of $M$ on input $x'$, equality of their configuration codes
implies $x=x'$.
\end{theorem}

\begin{theorem}[Equality of configuration codes]
\lean{Complexity.ConfigCodec.code_eq_iff}
\label{Complexity.ConfigCodec.code_eq_iff}
\leanok
\uses{Complexity.ConfigCodec.code, Turing.Cfg, Turing.Cfg.InWindow, Turing.NDTM.coreSum}
Let $M$ be a finite Turing machine with Boolean tape symbols, and let
$n,s\in\mathbb{N}$. Let $x$ be a Boolean input of length $n$, and let $c,d$ be
configurations of $M$ on input $x$, each with every work-tape head within distance $s$
of the origin and every work tape blank outside that range. Their Boolean codes are
equal if and only if their core configurations are equal and their output summaries are
equal.
\end{theorem}

\begin{theorem}[Validity CNF Characterizes Configuration Codes]
\lean{Complexity.ConfigCodec.validity_eval_iff}
\label{Complexity.ConfigCodec.validity_eval_iff}
\leanok
\uses{Complexity.ConfigCodec.code, Complexity.ConfigCodec.validity, Complexity.ConfigCodec.width, Turing.Cfg, Turing.Cfg.InWindow}
Let $M$ be a finite-state Turing machine with Boolean tape symbols, and let $n,s$ be
natural numbers. For every Boolean word $w$ of length equal to the codec width for
$M,n,s$, the validity CNF evaluates to true under the assignment that gives variable $v$
the bit of $w$ at position $v$ (and gives false when that position is outside $w$) if
and only if there are a Boolean input word $x$ of length $n$ and a configuration $c$ of
$M$ on input $x$ such that every work-tape head lies within distance $s$ of the origin,
every work-tape cell outside that interval is blank, and $w$ is the configuration code
of $x,c$ for parameters $n,s$.
\end{theorem}

\begin{theorem}[Adjacency formula characterizes a machine step]
\lean{Complexity.ConfigCodec.adjacency_eval_iff}
\label{Complexity.ConfigCodec.adjacency_eval_iff}
\leanok
\uses{Complexity.ConfigCodec.adjacency, Complexity.ConfigCodec.code, Turing.Cfg, Turing.Cfg.InWindow, Turing.MultiTapeTM.step, Turing.NDTM.coreSum}
Let $M$ be a finite Turing machine with Boolean tape symbols, and let $n,s$ be natural
numbers. Let $x$ be a Boolean input of length $n$, and let $c,d$ be configurations of
$M$ on input $x$, each with every work-tape head at distance at most $s$ from the origin
and every work tape blank at positions whose distance from the origin is greater than
$s$. The adjacency formula evaluated on the concatenation of the Boolean codes of $c$
and $d$ is true if and only if the core configuration and output summary of the one-step
successor of $c$ agree with those of $d$.
\end{theorem}

\begin{theorem}[Adjacency rejects distinct inputs]
\lean{Complexity.ConfigCodec.adjacency_eval_of_ne}
\label{Complexity.ConfigCodec.adjacency_eval_of_ne}
\leanok
\uses{Complexity.ConfigCodec.adjacency, Complexity.ConfigCodec.code, Turing.Cfg}
Let $M$ be a finite Turing machine with Boolean tape symbols, and let $n,s$ be natural
numbers. Let $x$ and $x'$ be distinct Boolean inputs, each of length $n$, and let $c$
and $c'$ be arbitrary configurations of $M$ on inputs $x$ and $x'$, respectively. The
adjacency formula evaluated on the concatenation of their Boolean codes is false.
\end{theorem}

\begin{theorem}[Acceptance condition for a configuration]
\lean{Complexity.ConfigCodec.acceptance_eval_iff}
\label{Complexity.ConfigCodec.acceptance_eval_iff}
\leanok
\uses{Complexity.ConfigCodec.acceptance, Complexity.ConfigCodec.code, Turing.Cfg, Turing.outSummary}
Let $M$ be a finite Turing machine with Boolean tape symbols, let $n,s\in\mathbb{N}$,
let $x$ be a finite Boolean input, and let $c$ be a configuration of $M$ on input $x$.
The acceptance formula evaluated on the Boolean code of $c$ is true if and only if $c$
is in the halting state and its output is exactly the one-symbol list $[\mathrm{true}]$.
\end{theorem}

\begin{theorem}[Variable Count Bound for the Validity Formula]
\lean{Complexity.ConfigCodec.validity_numVars_le}
\label{Complexity.ConfigCodec.validity_numVars_le}
\leanok
\uses{Complexity.ConfigCodec.validity, Complexity.ConfigCodec.width, Std.Sat.CNF.numVars}
For every finite-state Turing machine $M$ with Boolean tape symbols and all natural
numbers $n,s$, the variable count of the validity CNF for parameters $n,s$ is at most
the codec width for $M,n,s$.
\end{theorem}

\begin{theorem}[Variable bound for the acceptance formula]
\lean{Complexity.ConfigCodec.acceptance_numVars_le}
\label{Complexity.ConfigCodec.acceptance_numVars_le}
\leanok
\uses{Complexity.ConfigCodec.acceptance, Complexity.ConfigCodec.width, Std.Sat.CNF.numVars}
Let $M$ be a finite Turing machine with Boolean tape symbols, and let
$n,s\in\mathbb{N}$. The variable count of its acceptance formula is at most the width
$w(n,s)$, where $w(n,s)$ is the width assigned to encodings of configurations for these
parameters.
\end{theorem}

\begin{theorem}[Variable bound for the adjacency formula]
\lean{Complexity.ConfigCodec.adjacency_numVars_le}
\label{Complexity.ConfigCodec.adjacency_numVars_le}
\leanok
\uses{Complexity.ConfigCodec.adjacency, Complexity.ConfigCodec.width, Std.Sat.CNF.numVars}
Let $M$ be a finite-state Turing machine with Boolean tape symbols, and let
$n,s\in\mathbb{N}$. The variable count of the CNF formula encoding adjacency for
parameters $n$ and $s$ is at most twice the width of the configuration encoding,
$2\,\mathrm{width}(n,s)$.
\end{theorem}

\begin{theorem}[Polynomial length of serialized formulas]
\lean{Complexity.ConfigCodec.exists_serialize_length_le}
\label{Complexity.ConfigCodec.exists_serialize_length_le}
\leanok
\uses{Complexity.ConfigCodec.acceptance, Complexity.ConfigCodec.adjacency, Complexity.ConfigCodec.validity, Std.Sat.CNF.serialize}
For every finite-state Turing machine $M$ with Boolean tape symbols, there exists
$K\in\mathbb{N}$ such that for all $n,s\in\mathbb{N}$, the serialized validity,
adjacency, and acceptance formulas each have length at most $K(s+n+2)^K$.
\end{theorem}

\begin{theorem}[Polynomial-Time Encoding of Validity Constraints]
\lean{Complexity.ConfigCodec.polyTimeComputable_serialize_validity}
\label{Complexity.ConfigCodec.polyTimeComputable_serialize_validity}
\leanok
\uses{Complexity.ConfigCodec.validity, Complexity.PolyTimeComputable, Std.Sat.CNF.serialize}
Let $M$ be a finite-state Turing machine with Boolean tape symbols, and let
$n,s,o:\{0,1	ext{-strings}	o\mathbb N\}$ be functions such that each of the three maps
$x\mapsto 1^{n(x)}$, $x\mapsto 1^{s(x)}$, and $x\mapsto 1^{o(x)}$ is computable in
polynomial time in the stated machine model. Then the map sending $x$ to the binary
encoding of the validity CNF for parameters $n(x),s(x)$, with every variable index
shifted by $o(x)$, is polynomial-time computable.
\end{theorem}

\begin{theorem}[Polynomial-Time Adjacency Encoding]
\lean{Complexity.ConfigCodec.polyTimeComputable_serialize_adjacency}
\label{Complexity.ConfigCodec.polyTimeComputable_serialize_adjacency}
\leanok
\uses{Complexity.ConfigCodec.adjacency, Complexity.PolyTimeComputable, Std.Sat.CNF.serialize}
Let $M$ be a finite-state Turing machine with Boolean tape symbols, and let $n,s,o$
assign a natural number to each Boolean string. Suppose the functions $x\mapsto
1^{n(x)}$, $x\mapsto 1^{s(x)}$, and $x\mapsto 1^{o(x)}$ are each computable in
polynomial time by finite-state Turing machines. Then the function mapping $x$ to the
serialized CNF obtained from the adjacency constraints for parameters $n(x),s(x)$ by
shifting every variable index by $o(x)$ is computable in polynomial time.
\end{theorem}

\begin{theorem}[Polynomial-Time Encoding of Acceptance Constraints]
\lean{Complexity.ConfigCodec.polyTimeComputable_serialize_acceptance}
\label{Complexity.ConfigCodec.polyTimeComputable_serialize_acceptance}
\leanok
\uses{Complexity.ConfigCodec.acceptance, Complexity.PolyTimeComputable, Std.Sat.CNF.serialize}
Let $M$ be a finite-state Turing machine with Boolean tape symbols, and let
$n,s,o:\{0,1	ext{-strings}	o\mathbb N\}$ be functions such that each of the three maps
$x\mapsto 1^{n(x)}$, $x\mapsto 1^{s(x)}$, and $x\mapsto 1^{o(x)}$ is computable by
$M$-model polynomial time. Then the map sending $x$ to the binary encoding of the
acceptance CNF for parameters $n(x),s(x)$, with every variable index shifted by $o(x)$,
is polynomial-time computable.
\end{theorem}

\begin{theorem}[Polynomial time of initial configuration encoding]
\lean{Complexity.ConfigCodec.polyTimeComputable_code_initCfg}
\label{Complexity.ConfigCodec.polyTimeComputable_code_initCfg}
\leanok
\uses{Complexity.ConfigCodec.code, Complexity.PolyTimeComputable, Turing.MultiTapeTM.initCfg}
Let $M$ be a finite Turing machine over the Boolean alphabet, and let $s$ assign a
natural number to each Boolean string. If the function sending each string $x$ to the
string of $s(x)$ true bits is computable in polynomial time, then the function sending
$x$ to the configuration code of the initial configuration of $M$ on input $x$, using
parameters $n=|x|$ and space bound $s(x)$, is computable in polynomial time.
\end{theorem}

\begin{definition}[True Quantified Boolean Formula Language]
\lean{Complexity.TQBF}
\label{Complexity.TQBF}
\leanok
\uses{Complexity.QBF.decode, Complexity.QBF.truth}
The true quantified Boolean formula language is the set of Boolean strings whose decoded
quantified Boolean formula is true.
\end{definition}

\begin{theorem}[Polynomial-space decidability of quantified Boolean formulas]
\lean{Complexity.TQBF_mem_PSPACE}
\label{Complexity.TQBF_mem_PSPACE}
\leanok
\uses{Complexity.LogProg.TQBFARM.truth_mem_PSPACE, Complexity.PSPACE, Complexity.TQBF}
\difficulty{1}
The language of bit strings that decode to true quantified Boolean formulas is decidable
by a finite Turing machine using at most $c(n^c+1)$ work-tape cells on inputs of length
$n$, for some natural number $c$.
\end{theorem}

\begin{theorem}[TQBF is PSPACE-hard]
\lean{Complexity.TQBF_PSPACEHard}
\label{Complexity.TQBF_PSPACEHard}
\leanok
\uses{Complexity.PSPACEHard, Complexity.TQBF, Complexity.TQBFHard.PSPACE_polyTimeReducible_TQBF}
\difficulty{1}
The language of bit strings whose decoded quantified Boolean formula is true is
PSPACE-hard: every language in PSPACE has a polynomial-time many-one reduction to it.
\end{theorem}

\begin{theorem}[TQBF is PSPACE-complete]
\lean{Complexity.TQBF_PSPACEComplete}
\label{Complexity.TQBF_PSPACEComplete}
\leanok
\uses{Complexity.PSPACEComplete, Complexity.TQBF, Complexity.TQBF_PSPACEHard, Complexity.TQBF_mem_PSPACE}
\difficulty{1}
The language of bit strings whose decoded quantified Boolean formula is true is
PSPACE-complete: it belongs to PSPACE, and every language in PSPACE has a
polynomial-time many-one reduction to it.
\end{theorem}
